% =====================================================================
%  Dokumen Pengujian Perangkat Lunak (Software Testing Report)
%  ManganYuk! - Kelompok 6 - II4011 Transformasi Digital
%
%  Kompilasi:  pdflatex main.tex  (jalankan 2x agar daftar isi terisi)
%              atau latexmk -pdf main.tex
% =====================================================================
\documentclass[11pt,letterpaper,oneside]{report}

% ---------- Bahasa & font ----------
\usepackage[T1]{fontenc}
\usepackage[utf8]{inputenc}
\usepackage[indonesian]{babel}
\usepackage{lmodern}
\usepackage{amssymb}
\usepackage{textcomp}

% ---------- Tata letak ----------
\usepackage[letterpaper,margin=2.5cm]{geometry}
\usepackage{setspace}
\setstretch{1.15}
\setlength{\parindent}{0pt}
\setlength{\parskip}{6pt}

% ---------- Tabel ----------
\usepackage{array}
\usepackage{longtable}
\IfFileExists{colortbl.sty}{\usepackage[table]{xcolor}}{%
  \usepackage{xcolor}\newcommand{\cellcolor}[2][]{}\newcommand{\rowcolor}[2][]{}}
\usepackage{enumitem}
\newcolumntype{L}[1]{>{\raggedright\arraybackslash}p{\dimexpr#1\linewidth-2\tabcolsep-1.5\arrayrulewidth\relax}}
\newcolumntype{C}[1]{>{\centering\arraybackslash}p{\dimexpr#1\linewidth-2\tabcolsep-1.5\arrayrulewidth\relax}}
\renewcommand{\arraystretch}{1.25}
\setlist{nosep}
\setlist[itemize]{leftmargin=1.2em}
\AtBeginEnvironment{longtable}{\small\sloppy}

% ---------- Gambar ----------
\usepackage{graphicx}
\usepackage{float}
\usepackage{tikz}
\usetikzlibrary{shapes.geometric, positioning, fit, backgrounds, arrows.meta}
\graphicspath{{images/}}

% ---------- Judul bab, header & footer ----------
\usepackage{titlesec}
\titleformat{\chapter}[hang]{\Large\bfseries}{\thechapter.}{0.6em}{}
\titlespacing*{\chapter}{0pt}{0pt}{12pt}
\titleformat{\section}{\large\bfseries}{\thesection}{0.6em}{}
\titleformat{\subsection}{\normalsize\bfseries}{\thesubsection}{0.6em}{}

\usepackage{fancyhdr}
\pagestyle{fancy}
\fancyhf{}
\renewcommand{\headrulewidth}{0pt}
\fancyfoot[L]{\footnotesize\textit{Dokumen Pengujian Perangkat Lunak -- ManganYuk! (Kelompok 6)}}
\fancyfoot[R]{\footnotesize\thepage}
\fancypagestyle{plain}{\fancyhf{}\fancyfoot[L]{\footnotesize\textit{Dokumen Pengujian Perangkat Lunak -- ManganYuk! (Kelompok 6)}}\fancyfoot[R]{\footnotesize\thepage}}

% ---------- Caption & tautan ----------
\usepackage[font=small,labelfont=bf]{caption}
\usepackage[hidelinks]{hyperref}
\usepackage{etoolbox}

\hypersetup{
  pdftitle={Dokumen Pengujian Perangkat Lunak ManganYuk! - Kelompok 6},
  pdfsubject={II4011 Transformasi Digital},
}

% ---------- Simbol kotak centang untuk status test case ----------
\newcommand{\statusbox}{\(\square\) Pass \quad \(\square\) Fail \quad \(\square\) Blocked}
\newcommand{\passbox}{\(\boxtimes\) Pass \quad \(\square\) Fail \quad \(\square\) Blocked}
\newcommand{\statuscol}{\(\square\) Pass\newline \(\square\) Fail\newline \(\square\) Blocked}
\newcommand{\passcol}{\(\boxtimes\) Pass\newline \(\square\) Fail\newline \(\square\) Blocked}

% ---------- Halaman landscape (kertas diputar, isi tetap tegak) ----------
\makeatletter
\newlength{\PortraitPaperW}\newlength{\PortraitPaperH}
\newlength{\PortraitTextW}\newlength{\PortraitTextH}
\AtBeginDocument{%
  \setlength{\PortraitPaperW}{\paperwidth}\setlength{\PortraitPaperH}{\paperheight}%
  \setlength{\PortraitTextW}{\textwidth}\setlength{\PortraitTextH}{\textheight}}
\newcommand{\@setpagelayout}[4]{% paperW, paperH, textW, textH
  \setlength{\paperwidth}{#1}\setlength{\paperheight}{#2}%
  \ifdefined\pdfpagewidth \pdfpagewidth=#1 \pdfpageheight=#2 \fi
  \setlength{\textwidth}{#3}\setlength{\textheight}{#4}%
  \setlength{\hsize}{\textwidth}\setlength{\columnwidth}{\textwidth}%
  \setlength{\linewidth}{\textwidth}\setlength{\headwidth}{\textwidth}%
  \global\@colht\textheight\global\@colroom\textheight\global\vsize\textheight}
\newenvironment{landscapepage}{%
  \clearpage
  \@setpagelayout{\PortraitPaperH}{\PortraitPaperW}%
    {\dimexpr\PortraitTextW+\PortraitPaperH-\PortraitPaperW\relax}%
    {\dimexpr\PortraitTextH-\PortraitPaperH+\PortraitPaperW\relax}%
}{%
  \clearpage
  \@setpagelayout{\PortraitPaperW}{\PortraitPaperH}{\PortraitTextW}{\PortraitTextH}%
}
\makeatother

\begin{document}

% ===================== Bagian awal =====================
\pagenumbering{roman}
% Halaman sampul
\begin{titlepage}
\centering
\setstretch{1.2}

{\Large\bfseries DOKUMEN PENGUJIAN PERANGKAT LUNAK\par}
\vspace{0.2cm}
{\large\textit{(Software Testing Report)}\par}
\vspace{0.8cm}
{\Huge\bfseries ManganYuk!\par}
\vspace{0.3cm}
{\large Sudirman Street Day and Night Market\par}
\vspace{0.8cm}
Mata Kuliah: II4011 Transformasi Digital\par
Dosen Pengajar: Prof. Ir. Kridanto Surendro, M.Sc., Ph.D.\par

\vspace{1cm}
\includegraphics[width=4cm]{images/logo-itb.png}\par
\vspace{1cm}

\textbf{Disusun Oleh:}\par
[6 / ManganYuk!]\par
\vspace{0.3cm}
\renewcommand{\arraystretch}{1.3}
\begin{tabular}{|c|l|c|}
\hline
\rowcolor[HTML]{D9EAF7}\textbf{No.} & \textbf{Nama Anggota} & \textbf{NIM} \\ \hline
1 & Dita Maheswari             & 13523125 \\ \hline
2 & Samantha Laqueenna Ginting & 13523138 \\ \hline
3 & Anggita Najmi Layali       & 18223122 \\ \hline
4 & Theresia Ivana M. S        & 18223126 \\ \hline
5 & Raka Adhitya Nugraha       & 18223013 \\ \hline
\end{tabular}

\vfill
{\bfseries PROGRAM STUDI SISTEM DAN TEKNOLOGI INFORMASI\\
SEKOLAH TEKNIK ELEKTRO DAN INFORMATIKA\\
INSTITUT TEKNOLOGI BANDUNG\\
2026\par}
\end{titlepage}

% Daftar isi, daftar tabel, dan daftar gambar
\setcounter{tocdepth}{2}
\tableofcontents
\clearpage
\listoftables
\clearpage
\listoffigures


% ===================== Isi (per bab) =====================
\clearpage
\pagenumbering{arabic}
% Bab 1: Informasi Pengujian
\chapter{Informasi Pengujian}
\label{bab:informasi-pengujian}

\section{Tujuan dan Ruang Lingkup}

\begin{longtable}{|L{0.28}|L{0.72}|}
\caption{Tujuan dan Ruang Lingkup Pengujian}\\
\hline
\cellcolor[HTML]{D9EAF7}\textbf{Komponen} & \cellcolor[HTML]{D9EAF7}\textbf{Keterangan} \\ \hline
\endfirsthead
\hline
\cellcolor[HTML]{D9EAF7}\textbf{Komponen} & \cellcolor[HTML]{D9EAF7}\textbf{Keterangan} \\ \hline
\endhead
\textbf{Nama produk} & \textbf{ManganYuk!}, aplikasi \textit{location-based culinary \& local culture explorer} untuk kawasan Sudirman Street Food, Bandung. Aplikasi mobile Android (React Native + Expo). \\ \hline
\textbf{Versi / build yang diuji} & Branch \texttt{main} (Release v.x) \\ \hline
\textbf{Requirement dalam lingkup} &
\textbf{Requirement Fungsional (F01--F12):}
\begin{itemize}
  \item F01 Location Awareness
  \item F02 Location-Based Trigger (Geofence)
  \item F03 Multimodal Exploration
  \item F04 Eyes-Free Interaction
  \item F05 Search \& Destination
  \item F06 Navigation
  \item F07 Contextual Recommendation
  \item F08 Domain-Specific Innovative Feature: Smart Culinary Trail Planner
  \item F09 Location Privacy \& Consent
  \item F10 Paspor Kuliner \& Ulasan Tenant (Gamification)
  \item F11 Akun Pengguna Lokal
  \item F12 Preferensi Halal
\end{itemize} \\ \hline
\textbf{Di luar lingkup} &
\begin{itemize}
  \item Sistem pemesanan makanan online (\textit{online food ordering}) dan pembayaran digital (\textit{payment gateway}) (NG01).
  \item Fitur AI/LLM untuk sistem rekomendasi atau chatbot (NG02).
  \item Pembaruan stok menu atau antrean secara \textit{real-time} dari pihak tenant (NG03).
  \item Area kuliner lain di Kota Bandung selain \textit{pilot project} Sudirman Street Food (NG04).
  \item Reservasi tempat duduk atau meja secara langsung (NG05).
  \item Platform selain Android (iOS).
\end{itemize} \\ \hline
\textbf{Periode pengujian} & Sprint 5 (Week 14), 18 -- 25 November 2026 (Alpha Testing Android: 24 November 2026; Unit/Static Test: 25 November 2026; Beta Testing: menyusul) \\ \hline
\textbf{Lokasi pengujian lapangan} & Sudirman Street Day and Night Market, Bandung \\ \hline
\end{longtable}

\section{Metode Pengujian}

\begin{longtable}{|L{0.17}|L{0.17}|L{0.26}|L{0.20}|L{0.20}|}
\caption{Metode Pengujian}\\
\hline
\cellcolor[HTML]{D9EAF7}\textbf{Jenis} & \cellcolor[HTML]{D9EAF7}\textbf{Pelaksana} & \cellcolor[HTML]{D9EAF7}\textbf{Tujuan} & \cellcolor[HTML]{D9EAF7}\textbf{Teknik} & \cellcolor[HTML]{D9EAF7}\textbf{Output} \\ \hline
\endfirsthead
\hline
\cellcolor[HTML]{D9EAF7}\textbf{Jenis} & \cellcolor[HTML]{D9EAF7}\textbf{Pelaksana} & \cellcolor[HTML]{D9EAF7}\textbf{Tujuan} & \cellcolor[HTML]{D9EAF7}\textbf{Teknik} & \cellcolor[HTML]{D9EAF7}\textbf{Output} \\ \hline
\endhead
Alpha testing & Tim pengembang & Memastikan setiap fungsi berjalan sesuai \textit{acceptance criteria} & \textit{Black box}, berbasis \textit{test case} & Tabel hasil \textit{test case}, daftar \textit{defect} \\ \hline
Retest / regression & Tim pengembang & Memastikan perbaikan \textit{defect} berhasil dan tidak merusak fungsi lain & Menjalankan ulang \textit{test case} terkait & Hasil retest \\ \hline
Beta testing (UAT) & Calon pengguna (target user PRD) & Mengukur tingkat penerimaan pengguna & Tugas penggunaan + kuesioner skala Likert 1--5 & Indeks penerimaan per aspek \\ \hline
\end{longtable}

\section{Lingkungan Pengujian}

\begin{longtable}{|L{0.28}|L{0.72}|}
\caption{Lingkungan Pengujian}\\
\hline
\cellcolor[HTML]{D9EAF7}\textbf{Komponen} & \cellcolor[HTML]{D9EAF7}\textbf{Spesifikasi} \\ \hline
\endfirsthead
\hline
\cellcolor[HTML]{D9EAF7}\textbf{Komponen} & \cellcolor[HTML]{D9EAF7}\textbf{Spesifikasi} \\ \hline
\endhead
Perangkat uji & Android, GPS aktif \\ \hline
Versi aplikasi & Branch \texttt{main} (Release v.x); lingkungan uji otomatis: Node.js + Jest (frontend) \& Vitest (backend) \\ \hline
Backend / sumber data & Dataset lokal \texttt{frontend/src/data/pois.json} (31 POI dibundel di dalam APK); tile peta OpenStreetMap (MapLibre); layanan rute jalan kaki OSRM publik; serta backend Mangan Yuk (Supabase) untuk sinkronisasi Paspor Kuliner dan ulasan (F10) \\ \hline
Teknologi pemicu lokasi & GPS / geofencing (background \textit{foreground service}, akurasi tinggi) \\ \hline
Kondisi lapangan & Outdoor di area Sudirman Street Food, pukul 10.00 WIB (atau jam operasional tenant) dengan kondisi sinyal baik \\ \hline
\end{longtable}

\section{Tim Penguji}

\begin{longtable}{|C{0.08}|L{0.34}|C{0.18}|L{0.40}|}
\caption{Tim Penguji}\\
\hline
\cellcolor[HTML]{D9EAF7}\textbf{No} & \cellcolor[HTML]{D9EAF7}\textbf{Nama} & \cellcolor[HTML]{D9EAF7}\textbf{NIM} & \cellcolor[HTML]{D9EAF7}\textbf{Peran dalam Pengujian} \\ \hline
\endfirsthead
\hline
\cellcolor[HTML]{D9EAF7}\textbf{No} & \cellcolor[HTML]{D9EAF7}\textbf{Nama} & \cellcolor[HTML]{D9EAF7}\textbf{NIM} & \cellcolor[HTML]{D9EAF7}\textbf{Peran dalam Pengujian} \\ \hline
\endhead
1 & Dita Maheswari             & 13523125 & Alpha Testing \\ \hline
2 & Samantha Laqueenna Ginting & 13523138 & Alpha Testing \\ \hline
3 & Anggita Najmi Layali       & 18223122 & Alpha Testing \\ \hline
4 & Theresia Ivana M. S        & 18223126 & Alpha Testing \\ \hline
5 & Raka Adhitya Nugraha       & 18223013 & Alpha Testing \\ \hline
\end{longtable}

% Bab 2: Use Case
\chapter{Use Case}
\label{bab:use-case}

\section{Diagram Use Case}

Diagram use case berikut memodelkan interaksi antara aktor utama (\textbf{Pengunjung}) dan sistem eksternal (\textbf{Layanan Rute OSRM} serta \textbf{Backend Mangan Yuk}) terhadap 15 use case pada aplikasi \textbf{ManganYuk!}, konsisten dengan \textit{user flow} pada PRD Bagian 2.

\begin{figure}[H]
\centering
\resizebox{\textwidth}{!}{%
\begin{tikzpicture}[
  font=\sffamily\footnotesize,
  uc/.style={ellipse, draw=black!75, fill=blue!6, thick, minimum width=4.6cm, minimum height=0.82cm, align=center, inner sep=2pt},
  ucopt/.style={ellipse, draw=black!75, dashed, fill=teal!6, thick, minimum width=4.6cm, minimum height=0.82cm, align=center, inner sep=2pt},
  actor/.style={align=center},
  assoc/.style={-, thick, draw=black!65},
  ext/.style={->, dashed, thick, draw=black!55}
]

% Use cases - Kolom Kiri (Core Onboarding, Geofence, Eksplorasi)
\node[uc] (uc01) at (0, 11.2) {\textbf{UC-01} Onboarding \& Izin Lokasi};
\node[uc] (uc02) at (0, 9.9)  {\textbf{UC-02} Posisi \& Status Area};
\node[uc] (uc03) at (0, 8.6)  {\textbf{UC-03} Mode Jelajah \& Entry Event};
\node[uc] (uc04) at (0, 7.3)  {\textbf{UC-04} Kartu POI Proximity (10 m)};
\node[uc] (uc05) at (0, 6.0)  {\textbf{UC-05} Catat Kunjungan (Dwell/Exit)};
\node[uc] (uc06) at (0, 4.7)  {\textbf{UC-06} Kartu POI Multimodal};
\node[uc] (uc07) at (0, 3.4)  {\textbf{UC-07} Narasi Eyes-Free};
\node[uc] (uc12) at (0, 2.1)  {\textbf{UC-12} Matikan Mode Jelajah \& Privasi};

% Use cases - Kolom Kanan (Search, Nav, Rec, Trail, Optional)
\node[uc] (uc08) at (5.8, 11.2) {\textbf{UC-08} Cari Tenant \& Tujuan};
\node[uc] (uc09) at (5.8, 9.9)  {\textbf{UC-09} Navigasi Dalam Aplikasi};
\node[uc] (uc10) at (5.8, 8.6)  {\textbf{UC-10} Rekomendasi Berikutnya};
\node[uc] (uc11) at (5.8, 7.3)  {\textbf{UC-11} Smart Culinary Trail};
\node[ucopt] (uc13) at (5.8, 5.4) {\textbf{UC-13} Paspor Kuliner \& Ulasan (F10)};
\node[ucopt] (uc14) at (5.8, 4.1) {\textbf{UC-14} Akun Pengguna Lokal (F11)};
\node[ucopt] (uc15) at (5.8, 2.8) {\textbf{UC-15} Preferensi Halal (F12)};

% System Boundary
\begin{scope}[on background layer]
  \node[draw=black!80, thick, rounded corners=6pt, fill=gray!4,
        fit=(uc01)(uc12)(uc08)(uc15), inner sep=14pt,
        label={[font=\sffamily\bfseries\small, yshift=-14pt]north:Sistem Aplikasi ManganYuk! (Android)}] (sysbox) {};
\end{scope}

% Stick Figure Macro
\newcommand{\stickman}[3]{% name, x, y
  \begin{scope}[shift={(#2,#3)}]
    \draw[thick] (0,0.45) circle (0.18);
    \draw[thick] (0,0.27) -- (0,-0.25);
    \draw[thick] (-0.25,0.10) -- (0.25,0.10);
    \draw[thick] (0,-0.25) -- (-0.20,-0.62);
    \draw[thick] (0,-0.25) -- (0.20,-0.62);
    \node[below=0.70cm, align=center, font=\sffamily\bfseries\footnotesize] (#1) at (0,0) {#1};
  \end{scope}
}

% Aktor Kiri: Pengunjung
\begin{scope}[shift={(-4.6,6.8)}]
  \draw[thick] (0,0.45) circle (0.18);
  \draw[thick] (0,0.27) -- (0,-0.25);
  \draw[thick] (-0.25,0.10) -- (0.25,0.10);
  \draw[thick] (0,-0.25) -- (-0.20,-0.62);
  \draw[thick] (0,-0.25) -- (0.20,-0.62);
  \node[align=center, font=\sffamily\bfseries\footnotesize] at (0,-0.88) {Pengunjung};
  \coordinate (pengunjung) at (0.35, 0);
\end{scope}

% Aktor Kanan 1: Layanan Rute OSRM
\begin{scope}[shift={(10.4,9.9)}]
  \node[draw=black!75, thick, fill=orange!10, rounded corners=3pt, minimum width=2.4cm, minimum height=1.0cm, align=center, font=\sffamily\bfseries\footnotesize] (osrm) {<<system>>\\Layanan Rute\\OSRM};
\end{scope}

% Aktor Kanan 2: Backend Mangan Yuk
\begin{scope}[shift={(10.4,5.4)}]
  \node[draw=black!75, thick, fill=orange!10, rounded corners=3pt, minimum width=2.4cm, minimum height=1.0cm, align=center, font=\sffamily\bfseries\footnotesize] (backend) {<<system>>\\Backend\\Mangan Yuk};
\end{scope}

% Asosiasi Pengunjung -> UC kiri
\foreach \u in {uc01,uc02,uc03,uc04,uc05,uc06,uc07,uc12} {
  \draw[assoc] (pengunjung) -- (\u.west);
}

% Asosiasi Pengunjung -> UC kanan (jalur rapi via sisi atas/bawah)
\foreach \u in {uc08,uc09,uc10,uc11,uc13,uc14,uc15} {
  \draw[assoc] (pengunjung) -- (\u.west);
}

% Asosiasi Aktor Eksternal
\draw[assoc] (uc09.east) -- (osrm.west);
\draw[assoc] (uc13.east) -- (backend.west);

% Legenda singkat di bawah
\node[anchor=north west, font=\sffamily\scriptsize] at ([yshift=-6pt]sysbox.south west)
  {Keterangan: Garis solid = Requirement Utama (F01--F09); Garis putus-putus = Requirement Tambahan/Opsional (F10--F12).};

\end{tikzpicture}%
}
\caption{Diagram Use Case Aplikasi ManganYuk!}
\label{fig:use-case-diagram}
\end{figure}

\section{Daftar Use Case}

\begin{longtable}{|C{0.10}|L{0.25}|L{0.14}|L{0.13}|L{0.38}|}
\caption{Daftar Use Case}\\
\hline
\cellcolor[HTML]{D9EAF7}\textbf{UC ID} & \cellcolor[HTML]{D9EAF7}\textbf{Nama Use Case} & \cellcolor[HTML]{D9EAF7}\textbf{Aktor} & \cellcolor[HTML]{D9EAF7}\textbf{Related Req.} & \cellcolor[HTML]{D9EAF7}\textbf{Deskripsi Singkat} \\ \hline
\endfirsthead
\hline
\cellcolor[HTML]{D9EAF7}\textbf{UC ID} & \cellcolor[HTML]{D9EAF7}\textbf{Nama Use Case} & \cellcolor[HTML]{D9EAF7}\textbf{Aktor} & \cellcolor[HTML]{D9EAF7}\textbf{Related Req.} & \cellcolor[HTML]{D9EAF7}\textbf{Deskripsi Singkat} \\ \hline
\endhead
UC-01 & Melakukan onboarding dan memberi izin lokasi/notifikasi & Pengunjung & F09, F01, F12 & Pengguna membaca penjelasan privasi lokasi, memberi izin lokasi dan notifikasi, lalu memilih preferensi rasa, jenis kuliner, serta opsi halal. \\ \hline
UC-02 & Melihat posisi dan status di dalam/di luar area & Pengunjung & F01 & Aplikasi menampilkan posisi pengguna di peta MapLibre beserta status di dalam/di luar area dan jarak ke Sudirman Street Food. \\ \hline
UC-03 & Mengaktifkan Mode Jelajah dan menerima notifikasi entry & Pengunjung & F02, F09 & Pengguna mengaktifkan Mode Jelajah dan menerima notifikasi sambutan saat memasuki radius 100 m kawasan. \\ \hline
UC-04 & Menerima kartu POI saat mendekati tenant & Pengunjung & F02, F12 & Aplikasi menampilkan kartu POI dan tawaran audio saat pengguna memasuki radius 10 m tenant, tanpa terpicu berulang. \\ \hline
UC-05 & Mencatat kunjungan otomatis (dwell dan exit event) & Pengunjung & F02 & Aplikasi mencatat kunjungan otomatis bila pengguna berada \(\ge\)3 menit di geofence tenant, dan mencatat exit event saat keluar (> radius + 10 m). \\ \hline
UC-06 & Melihat kartu POI multimodal & Pengunjung & F03 & Pengguna membuka kartu POI berisi foto (jika tersedia), cerita budaya, menu \& harga, jam buka, serta audio narasi dan transkrip. \\ \hline
UC-07 & Mendengarkan narasi tanpa melihat layar (eyes-free) & Pengunjung & F04 & Pengguna mendengarkan prompt dan narasi tenant, serta panduan audio, tanpa harus membuka atau melihat layar. \\ \hline
UC-08 & Mencari tenant dan menetapkan tujuan & Pengunjung & F05, F12 & Pengguna mencari tenant berdasarkan nama/menu/kategori (toleran typo), memakai filter (termasuk opsi halal), dan menetapkan tenant sebagai tujuan. \\ \hline
UC-09 & Navigasi menuju kawasan dan tenant & Pengunjung, Layanan Rute OSRM & F06, F01 & Pengguna dipandu di dalam aplikasi menuju tenant (peta MapLibre, rute jalan kaki OSRM atau fallback garis lurus, jarak, panah kompas, audio), termasuk dari luar kawasan. \\ \hline
UC-10 & Menerima rekomendasi tenant berikutnya & Pengunjung & F07, F12 & Setelah keluar dari tenant yang dikunjungi, aplikasi menampilkan 3 pilihan tenant berikutnya beserta alasannya. \\ \hline
UC-11 & Membuat dan menjalankan Smart Culinary Trail & Pengunjung & F08, F02, F12 & Pengguna mengisi durasi, budget, dan rasa; aplikasi menyusun itinerary trail dan memantau progres secara otomatis. \\ \hline
UC-12 & Mematikan Mode Jelajah dan menjaga privasi lokasi & Pengunjung & F09 & Pengguna dapat mematikan Mode Jelajah kapan saja; pelacakan berhenti dan koordinat tidak dikirim ke backend Mangan Yuk. \\ \hline
UC-13 & Mengumpulkan stempel Paspor Kuliner dan memberi ulasan & Pengunjung, Backend Mangan Yuk & F10 & Pengguna memberi ulasan tenant (rating, menu, teks/foto) untuk memperoleh stempel digital dan membuka badge trail tematik. \\ \hline
UC-14 & Mengelola akun pengguna lokal & Pengunjung & F11 & Pengguna masuk atau mendaftar akun lokal, menyimpan sesi di perangkat, mengubah preferensi pada Profil, dan keluar akun. \\ \hline
UC-15 & Mengatur preferensi tenant halal & Pengunjung & F12 & Pengguna mengaktifkan atau menonaktifkan opsi ``Hanya tenant halal'' agar peta, pencarian, trail, rekomendasi, dan trigger menyaring tenant non-halal. \\ \hline
\end{longtable}

\section{Spesifikasi Use Case}

\newcommand{\ucspec}[9]{%
\begin{longtable}{|L{0.26}|L{0.74}|}
\caption{Spesifikasi Use Case #1}\\
\hline
\cellcolor[HTML]{D9EAF7}\textbf{Atribut} & \cellcolor[HTML]{D9EAF7}\textbf{Detail} \\ \hline
\endfirsthead
\hline
\cellcolor[HTML]{D9EAF7}\textbf{Atribut} & \cellcolor[HTML]{D9EAF7}\textbf{Detail} \\ \hline
\endhead
\textbf{Use Case ID} & #1 \\ \hline
\textbf{Nama Use Case} & #2 \\ \hline
\textbf{Aktor} & #3 \\ \hline
\textbf{Related Requirement} & #4 \\ \hline
\textbf{Deskripsi} & #5 \\ \hline
\textbf{Pre-condition} & #6 \\ \hline
\textbf{Trigger} & #7 \\ \hline
\textbf{Alur utama} & #8 \\ \hline
\textbf{Alur alternatif / eksepsi} & #9 \\ \hline
\end{longtable}
}

% UC-01
\ucspec{UC-01}
{Melakukan onboarding dan memberi izin lokasi/notifikasi}
{Pengunjung}
{F09 Location Privacy \& Consent; F01 Location Awareness; F12 Preferensi Halal}
{Pengguna membaca penjelasan privasi lokasi, memberi izin lokasi dan notifikasi, lalu memilih preferensi kuliner dan opsi halal.}
{Aplikasi baru dipasang (data aplikasi bersih); pengguna telah melewati layar login/registrasi (F11); izin lokasi dan notifikasi belum diberikan.}
{Pengguna masuk ke halaman onboarding untuk pertama kali.}
{1. Sistem menampilkan halaman onboarding (jelajah kuliner, cerita budaya, culinary trail).\newline
2. Pengguna menekan tombol Lanjut.\newline
3. Sistem menampilkan penjelasan penggunaan lokasi (alasan, kapan lokasi dibaca, riwayat lokasi tidak dikirim ke backend Mangan Yuk, koordinat hanya dipakai layanan rute OSRM saat navigasi) dan tombol ``Izinkan Lokasi''.\newline
4. Pengguna menekan ``Izinkan Lokasi'', lalu memilih ``Izinkan'' pada dialog sistem.\newline
5. Sistem meminta izin notifikasi; pengguna mengizinkan.\newline
6. Sistem meminta preferensi rasa, jenis kuliner, dan opsi ``Hanya tenant halal'' (dengan opsi Lewati); pengguna memilih lalu menekan Mulai.\newline
7. Sistem menyimpan preferensi secara lokal dan menampilkan Beranda.}
{4a. Pengguna menekan ``Tolak'' pada dialog izin lokasi: sistem menampilkan penjelasan bahwa trigger otomatis dan navigasi tidak tersedia, dengan tombol Buka Pengaturan dan Lanjut tanpa Lokasi.\newline
4b. Pengguna menekan ``Lanjut tanpa Lokasi'': sistem menampilkan peta tanpa posisi pengguna dan daftar tenant; kartu POI dan audio dapat dibuka manual.\newline
6a. Pengguna menekan Lewati pada preferensi: sistem melanjutkan ke Beranda.\newline\textbf{Post-condition:} Preferensi tersimpan secara lokal; Beranda tampil (mode dengan lokasi atau mode tanpa lokasi).}

% UC-02
\ucspec{UC-02}
{Melihat posisi dan status di dalam/di luar area}
{Pengunjung}
{F01 Location Awareness}
{Aplikasi menampilkan posisi pengguna di peta MapLibre beserta status di dalam/di luar area dan jarak ke Sudirman Street Food.}
{Aplikasi terbuka di Beranda; izin lokasi diberikan; GPS aktif.}
{Pengguna membuka Beranda atau berpindah posisi.}
{1. Sistem membaca posisi pengguna (pembaruan tiap \(\pm\)3 detik atau perpindahan \(\ge\)5 m).\newline
2. Sistem menampilkan titik biru dengan lingkaran akurasi pada peta MapLibre (tile OpenStreetMap) dan denah kawasan.\newline
3. Sistem menampilkan status ``Di dalam/Di luar area Sudirman Street Food''.\newline
4. Sistem memperbarui posisi dan status saat pengguna berjalan.}
{3a. Pengguna berada di luar kawasan: sistem menampilkan status ``Di luar area'' dan jarak ke Sudirman Street Food; pengguna tetap dapat memilih tenant untuk memulai navigasi rute OSRM dari luar kawasan (lihat UC-09).\newline\textbf{Post-condition:} Posisi dan status area tampil sesuai kondisi aktual pengguna.}

% UC-03
\ucspec{UC-03}
{Mengaktifkan Mode Jelajah dan menerima notifikasi entry}
{Pengunjung}
{F02 Location-Based Trigger (Geofence); F09 Location Privacy \& Consent}
{Pengguna mengaktifkan Mode Jelajah dan menerima notifikasi sambutan saat memasuki radius 100 m kawasan.}
{Onboarding selesai; izin lokasi dan notifikasi diberikan; pengguna menuju kawasan.}
{Pengguna mengaktifkan Mode Jelajah dan memasuki radius 100 m kawasan.}
{1. Pengguna mengaktifkan Mode Jelajah.\newline
2. Sistem mengaktifkan pelacakan lokasi di background (\textit{foreground service}) dan menampilkan notifikasi permanen ``Mangan Yuk! sedang memandu''.\newline
3. Pengguna memasuki radius 100 m kawasan.\newline
4. Sistem memeriksa akurasi pembacaan lokasi (\(\le\)30 m).\newline
5. Sistem menampilkan notifikasi entry dan audio sambutan (``Selamat datang di Sudirman Street Food!'').\newline
6. Sistem menawarkan pilihan Rencanakan Trail atau Jelajah Bebas.}
{4a. Akurasi >30 m: sistem mengabaikan pembacaan, menampilkan banner ``Sinyal GPS lemah, trigger otomatis ditunda''; pengguna dapat membuka kartu POI manual dari daftar ``Tenant di sekitar''; trigger otomatis berlanjut setelah akurasi \(\le\)30 m.\newline
6a. Pengguna memilih Jelajah Bebas: peta menampilkan seluruh marker tenant sesuai preferensi halal tanpa tujuan tertentu dan proximity event tetap aktif.\newline\textbf{Post-condition:} Mode Jelajah aktif; kedatangan di kawasan terdeteksi.}

% UC-04
\ucspec{UC-04}
{Menerima kartu POI saat mendekati tenant}
{Pengunjung}
{F02 Location-Based Trigger (Geofence); F12 Preferensi Halal}
{Aplikasi menampilkan kartu POI dan tawaran audio saat pengguna memasuki radius 10 m tenant, tanpa terpicu berulang.}
{Mode Jelajah aktif; izin lokasi diberikan; data 31 POI termuat dari \texttt{pois.json}.}
{Posisi pengguna masuk ke radius proximity tenant (10 m).}
{1. Sistem membaca posisi pengguna dari pembaruan lokasi background.\newline
2. Sistem menghitung jarak ke setiap POI (haversine) dan mengenali POI terdekat.\newline
3. Sistem memeriksa pengguna berada dalam radius 10 m, akurasi \(\le\)30 m, preferensi halal terpenuhi (jika F12 aktif), dan POI belum terpicu/tidak dalam cooldown.\newline
4. Sistem menampilkan kartu POI, getar singkat, dan tawaran audio (``Kamu mendekati [tenant]. Mau dengar ceritanya?'').\newline
5. Pengguna membuka detail tenant atau mendengarkan narasi.}
{3a. POI sudah pernah terpicu dan pengguna masih di area yang sama, atau POI dalam cooldown 30 menit: sistem tidak menampilkan notifikasi ulang.\newline
3b. Akurasi lokasi >30 m: pembacaan diabaikan dan trigger ditunda (lihat UC-03).\newline
3c. Opsi ``Hanya tenant halal'' aktif dan POI berstatus non-halal/belum diverifikasi: proximity event tidak dipicu.\newline\textbf{Post-condition:} POI tercatat sebagai sudah terpicu (cooldown 30 menit); kartu POI tampil.}

% UC-05
\ucspec{UC-05}
{Mencatat kunjungan otomatis (dwell dan exit event)}
{Pengunjung}
{F02 Location-Based Trigger (Geofence)}
{Aplikasi mencatat kunjungan otomatis bila pengguna berada \(\ge\)3 menit di geofence tenant dan mencatat exit event saat keluar.}
{Mode Jelajah aktif; pengguna berada di dalam geofence tenant (radius 10 m).}
{Pengguna berada di dalam geofence tenant.}
{1. Sistem menghitung durasi pengguna di dalam geofence tenant.\newline
2. Setelah \(\ge\)3 menit, sistem mencatat kunjungan (\textit{dwell event}) di perangkat dan menandai tenant sudah dikunjungi.\newline
3. Bila trail aktif, sistem menandai tenant pada trail sebagai selesai.\newline
4. Saat pengguna keluar dari geofence (jarak > radius + 10 m, yakni >20 m), sistem mencatat \textit{exit event}.}
{2a. Pengguna keluar dari geofence sebelum 3 menit (hanya lewat): kunjungan tidak tercatat.\newline\textbf{Post-condition:} Kunjungan tercatat pada riwayat lokal perangkat.}

% UC-06
\ucspec{UC-06}
{Melihat kartu POI multimodal}
{Pengunjung}
{F03 Multimodal Exploration}
{Pengguna membuka kartu POI berisi foto (jika tersedia), cerita budaya, menu \& harga, jam buka, serta audio narasi dan transkrip.}
{Aplikasi terbuka; data 31 POI dibundel di dalam APK.}
{Pengguna memilih tenant dari peta/daftar atau kartu POI muncul dari trigger.}
{1. Pengguna memilih tenant dari peta/daftar atau menerima kartu POI dari trigger.\newline
2. Sistem membuka kartu POI (foto makanan/tenant jika tersedia, kategori dan profil rasa, menu \& harga, jam buka).\newline
3. Pengguna memutar audio narasi TTS (45--90 detik) yang memuat cerita kuliner/budaya; audio dapat dijeda dan dilanjutkan.\newline
4. Pengguna membuka transkrip audio.}
{3a. Tidak ada koneksi internet: seluruh konten dataset, foto lokal, dan audio tetap dapat diakses karena dibundel di APK.\newline\textbf{Post-condition:} Informasi tenant ditampilkan secara lengkap.}

% UC-07
\ucspec{UC-07}
{Mendengarkan narasi tanpa melihat layar (eyes-free)}
{Pengunjung}
{F04 Eyes-Free Interaction}
{Pengguna mendengarkan prompt dan narasi tenant, serta panduan audio, tanpa harus membuka atau melihat layar.}
{Mode Jelajah aktif; izin lokasi dan notifikasi diberikan; layar terkunci atau HP di saku.}
{Proximity event pada tenant tujuan atau tenant yang dilewati.}
{1. Sistem memutar prompt suara (``Kamu mendekati [tenant]. Mau dengar ceritanya?'') dan HP bergetar.\newline
2. Notifikasi menampilkan tombol aksi Dengar/Lewati, termasuk pada lock screen.\newline
3. Pengguna menekan Dengar.\newline
4. Sistem memutar audio narasi tanpa membuka aplikasi.\newline
5. Pada Mode Jalan, narasi diputar otomatis; navigasi memberi panduan audio tiap pengurangan jarak 20 m dan saat tiba.}
{3a. Pengguna menekan Lewati: prompt ditutup, tenant masuk daftar ``Terlewati'', dan tidak dipicu ulang selama cooldown 30 menit.\newline\textbf{Post-condition:} Narasi terdengar (atau dilewati) tanpa pengguna membuka aplikasi.}

% UC-08
\ucspec{UC-08}
{Mencari tenant dan menetapkan tujuan}
{Pengunjung}
{F05 Search \& Destination; F12 Preferensi Halal}
{Pengguna mencari tenant berdasarkan nama/menu/kategori (toleran typo), memakai filter, dan menetapkan tenant sebagai tujuan.}
{Aplikasi terbuka; data 31 POI tersedia secara offline.}
{Pengguna menekan kolom pencarian.}
{1. Pengguna menekan kolom pencarian dan mengetik kata kunci.\newline
2. Sistem menampilkan hasil pencarian toleran typo, diurutkan berdasarkan jarak.\newline
3. Pengguna menerapkan filter (rentang harga, rasa, jenis, status buka).\newline
4. Sistem memperbarui daftar sesuai kombinasi filter (dan menyembunyikan tenant non-halal jika F12 aktif).\newline
5. Pengguna memilih tenant dan menekan tombol Jadikan Tujuan.\newline
6. Sistem masuk ke mode navigasi menuju tenant yang dipilih.}
{2a. Tidak ada tenant yang sesuai dengan kata kunci/filter: sistem menampilkan daftar kosong/pesan tidak ditemukan tanpa crash.\newline\textbf{Post-condition:} Tenant tujuan ditetapkan dan mode navigasi aktif.}

% UC-09
\ucspec{UC-09}
{Navigasi menuju kawasan dan tenant}
{Pengunjung, Layanan Rute OSRM}
{F06 Navigation; F01 Location Awareness}
{Pengguna dipandu di dalam aplikasi menuju tenant tujuan (peta MapLibre, rute jalan kaki OSRM atau fallback garis lurus, jarak, panah kompas, audio), termasuk saat memulai dari luar kawasan.}
{Izin lokasi diberikan; tujuan ditetapkan.}
{Pengguna menekan Jadikan Tujuan atau Mulai Trail.}
{1. Pengguna menetapkan tenant tujuan (dari dalam maupun luar kawasan).\newline
2. Sistem menampilkan rute jalan kaki OSRM di peta MapLibre dari posisi pengguna ke tenant, jarak dalam meter, dan panah arah berbasis kompas.\newline
3. Jika pengguna keluar >25 m dari rute, sistem menghitung ulang rute OSRM.\newline
4. Sistem memberi panduan audio setiap jarak berkurang 20 m.\newline
5. Saat jarak \(\le\)10 m, sistem memutar ``Kamu sudah sampai!'' dan membuka kartu POI tenant tujuan.}
{2a. Layanan OSRM gagal atau tidak ada koneksi internet: sistem menampilkan garis lurus ke tenant tujuan dan navigasi tetap berjalan.\newline
1a. Izin lokasi ditolak: navigasi tidak aktif (mode tanpa lokasi).\newline\textbf{Post-condition:} Pengguna tiba di tenant tujuan dan kartu POI terbuka.}

% UC-10
\ucspec{UC-10}
{Menerima rekomendasi tenant berikutnya}
{Pengunjung}
{F07 Contextual Recommendation; F12 Preferensi Halal}
{Setelah keluar dari tenant yang dikunjungi, aplikasi menampilkan 3 pilihan tenant berikutnya beserta alasannya.}
{Mode Jelajah aktif; minimal satu tenant sudah dikunjungi (dwell tercatat); preferensi tersimpan.}
{Exit event dari geofence tenant yang sudah dikunjungi (jarak >20 m).}
{1. Pengguna berjalan keluar dari geofence tenant yang sudah dikunjungi.\newline
2. Sistem menghitung skor \textit{rule-based} (jarak, status buka, kecocokan rasa, sisa budget trail, preferensi halal, status kunjungan, aturan pasangan rasa).\newline
3. Sistem menampilkan ``3 pilihan berikutnya'' beserta alasan masing-masing.\newline
4. Pengguna memilih tenant dan masuk ke mode navigasi.}
{2a. Tenant tutup, sudah dikunjungi, atau non-halal saat F12 aktif: tenant dikeluarkan dari kandidat rekomendasi.\newline\textbf{Post-condition:} Rekomendasi tampil; pengguna dapat memilih tujuan berikutnya.}

% UC-11
\ucspec{UC-11}
{Membuat dan menjalankan Smart Culinary Trail}
{Pengunjung}
{F08 Smart Culinary Trail Planner; F02 Location-Based Trigger (Geofence); F12 Preferensi Halal}
{Pengguna mengisi durasi, budget, dan rasa; aplikasi menyusun itinerary trail dan memantau progres secara otomatis.}
{Mode Jelajah aktif; pengguna berada di dalam kawasan; preferensi tersimpan.}
{Pengguna memilih Rencanakan Trail.}
{1. Pengguna memilih Rencanakan Trail (dari notifikasi entry atau Beranda).\newline
2. Sistem menampilkan form trail: durasi (30/60/90 menit), budget, preferensi rasa (terisi otomatis dari onboarding).\newline
3. Pengguna memilih durasi 60 menit, budget Rp50.000, lalu menekan Buat Trail.\newline
4. Sistem menampilkan itinerary tenant berurutan (\textit{nearest neighbor}, dengan status halal jika F12 aktif) beserta estimasi total biaya, estimasi waktu, dan rute di peta.\newline
5. Pengguna menekan Mulai Trail; sistem masuk mode navigasi tenant pertama.\newline
6. Setiap dwell event menandai tenant selesai; sistem mengirim pengingat sisa waktu bila dwell melebihi alokasi + 5 menit; di akhir ditampilkan ringkasan trail.}
{6a. Tenant tutup atau pengguna menekan Lewati Tenant: sistem menandai tenant dilewati dan menyusun ulang sisa trail dari posisi pengguna sesuai sisa waktu dan sisa budget (tenant yang sudah selesai/dilewati tidak dipilih lagi).\newline\textbf{Post-condition:} Progres trail tersimpan; ringkasan trail (jumlah tenant, estimasi belanja, durasi, jarak) tampil.}

% UC-12
\ucspec{UC-12}
{Mematikan Mode Jelajah dan menjaga privasi lokasi}
{Pengunjung}
{F09 Location Privacy \& Consent}
{Pengguna dapat mematikan Mode Jelajah kapan saja; pelacakan berhenti dan koordinat tidak dikirim ke backend Mangan Yuk.}
{Mode Jelajah aktif.}
{Pengguna menekan Selesai Jelajah.}
{1. Pengguna menekan tombol Selesai Jelajah.\newline
2. Sistem menghentikan pelacakan lokasi di background.\newline
3. Sistem menghapus notifikasi permanen.\newline
4. Sistem kembali ke Beranda.}
{Selama aplikasi digunakan, tidak ada request ke backend Mangan Yuk yang membawa koordinat pengguna; koordinat hanya dikirim ke layanan rute OSRM saat layar navigasi aktif.\newline\textbf{Post-condition:} Pelacakan background berhenti dan tidak ada koordinat pengguna yang dikirim ke backend Mangan Yuk.}

% UC-13
\ucspec{UC-13}
{Mengumpulkan stempel Paspor Kuliner dan memberi ulasan}
{Pengunjung, Backend Mangan Yuk}
{F10 Paspor Kuliner \& Ulasan Tenant (Gamification)}
{Pengguna memberi ulasan tenant untuk memperoleh stempel digital pada trail tematik dan membuka badge.}
{Pengguna sudah masuk akun lokal (F11); halaman Paspor Kuliner atau detail tenant terbuka.}
{Pengguna menekan Tambah Ulasan.}
{1. Pengguna membuka Paspor Kuliner dari Profil atau menekan Tambah Ulasan pada detail tenant.\newline
2. Pengguna memilih tenant, memberi rating 1--5, memilih menu yang dimakan (milik tenant tersebut), serta mengisi teks ulasan dan foto (opsional).\newline
3. Sistem menyimpan ulasan, mencatat stempel digital untuk tenant tersebut, dan memperbarui progres trail tematik (``Jejak Kuliner Nusantara'' atau ``Manis-manis Bandung'').\newline
4. Sistem menyinkronkan data ke backend Mangan Yuk (\texttt{/v1/passport} dan \texttt{/v1/reviews}).\newline
5. Saat seluruh 5 tenant pada satu trail berstempel, sistem membuka badge trail tersebut.}
{3a. Pengguna memberi ulasan kedua pada tenant yang sama: ulasan baru tersimpan tanpa menambah stempel ganda.\newline
4a. Backend tidak dapat dihubungi: ulasan dan stempel tetap tersimpan di perangkat (AsyncStorage).\newline\textbf{Post-condition:} Ulasan, stempel, dan progres trail tersimpan dan bertahan setelah aplikasi ditutup ulang.}

% UC-14
\ucspec{UC-14}
{Mengelola akun pengguna lokal}
{Pengunjung}
{F11 Akun Pengguna Lokal}
{Pengguna mendaftar atau masuk ke akun lokal, menyimpan sesi secara lokal di perangkat, dan mengelola profil.}
{Aplikasi terpasang di perangkat.}
{Pengguna membuka aplikasi tanpa sesi aktif atau membuka halaman Profil.}
{1. Sistem mengarahkan pengguna yang belum masuk ke halaman login.\newline
2. Pengguna mendaftar atau masuk dengan email, nama pengguna, dan kata sandi (opsional mengaktifkan ``Ingat Saya'').\newline
3. Sistem menyimpan sesi di perangkat melalui SecureStore (tanpa mengirim ke server) dan mengarahkan ke onboarding atau Beranda.\newline
4. Pada halaman Profil, pengguna dapat melihat data akun, ringkasan paspor, mengubah preferensi, atau menekan Keluar.}
{4a. Pengguna menekan Keluar: sistem menghapus sesi dari perangkat dan mengarahkan ke halaman login.\newline\textbf{Post-condition:} Sesi akun lokal tersimpan atau terhapus sesuai aksi pengguna.}

% UC-15
\ucspec{UC-15}
{Mengatur preferensi tenant halal}
{Pengunjung}
{F12 Preferensi Halal}
{Pengguna mengaktifkan atau menonaktifkan opsi ``Hanya tenant halal'' dari onboarding atau Profil.}
{Pengguna berada di halaman onboarding atau Profil.}
{Pengguna mengubah toggle ``Hanya tenant halal''.}
{1. Pengguna mengaktifkan opsi ``Hanya tenant halal''.\newline
2. Sistem menyaring marker dan carousel peta, hasil pencarian (F05), pilihan Smart Culinary Trail (F08), rekomendasi (F07), serta proximity event (F02) agar hanya memproses tenant berstatus halal.\newline
3. Pengguna dapat menonaktifkan kembali opsi tersebut dari halaman Profil untuk menampilkan seluruh tenant.}
{1a. Tenant berstatus non-halal atau belum diverifikasi: disembunyikan dari peta, pencarian, trail, dan rekomendasi serta tidak memicu notifikasi proximity saat opsi aktif.\newline\textbf{Post-condition:} Seluruh fitur terkait mematuhi status preferensi halal pengguna.}

% Bab 3: Skenario Pengujian
\chapter{Skenario Pengujian}
\label{bab:skenario-pengujian}

\begin{longtable}{|C{0.09}|C{0.09}|L{0.48}|C{0.11}|C{0.13}|C{0.10}|}
\caption{Daftar Skenario Pengujian}\\
\hline
\cellcolor[HTML]{D9EAF7}\textbf{SC ID} & \cellcolor[HTML]{D9EAF7}\textbf{UC ID} & \cellcolor[HTML]{D9EAF7}\textbf{Skenario Pengujian} & \cellcolor[HTML]{D9EAF7}\textbf{Tipe} & \cellcolor[HTML]{D9EAF7}\textbf{Teknik} & \cellcolor[HTML]{D9EAF7}\textbf{Jumlah TC} \\ \hline
\endfirsthead
\hline
\cellcolor[HTML]{D9EAF7}\textbf{SC ID} & \cellcolor[HTML]{D9EAF7}\textbf{UC ID} & \cellcolor[HTML]{D9EAF7}\textbf{Skenario Pengujian} & \cellcolor[HTML]{D9EAF7}\textbf{Tipe} & \cellcolor[HTML]{D9EAF7}\textbf{Teknik} & \cellcolor[HTML]{D9EAF7}\textbf{Jumlah TC} \\ \hline
\endhead
SC-01 & UC-01 & Pengguna menyelesaikan onboarding; penjelasan privasi (termasuk layanan rute OSRM) tampil sebelum dialog izin sistem; izin lokasi dan notifikasi diberikan & Positif & Black box & 1 \\ \hline
SC-02 & UC-01 & Pengguna menolak izin lokasi & Negatif & Black box & 1 \\ \hline
SC-03 & UC-02 & Posisi pengguna tampil di peta MapLibre dan diperbarui saat berjalan & Positif & Black box & 2 \\ \hline
SC-04 & UC-02 & Status di dalam/di luar area berubah saat pengguna melewati batas kawasan & Positif & Black box & 1 \\ \hline
SC-05 & UC-02 & Pengguna berada di luar kawasan: status ``Di luar area'' dan tidak ada notifikasi entry & Negatif & Black box & 1 \\ \hline
SC-06 & UC-03 & Mode Jelajah aktif dan notifikasi entry muncul saat pengguna memasuki radius 100 m (termasuk saat background/layar terkunci) & Positif & Black box & 1 \\ \hline
SC-07 & UC-03 & Pembacaan lokasi dengan akurasi >30 m diabaikan & Negatif & Black box & 1 \\ \hline
SC-08 & UC-04 & Pengguna memasuki radius 10 m POI dan kartu POI tampil & Positif & Black box & 1 \\ \hline
SC-09 & UC-04 & Pengguna tetap berada di area POI; trigger tidak muncul berulang & Negatif & Black box & 1 \\ \hline
SC-10 & UC-05 & Pengguna berada \(\ge\)3 menit di tenant dan kunjungan tercatat & Positif & Black box & 1 \\ \hline
SC-11 & UC-05 & Pengguna hanya lewat (<3 menit) sehingga kunjungan tidak tercatat & Negatif & Black box & 1 \\ \hline
SC-12 & UC-06 & Kartu POI menampilkan konten lengkap dan terbuka cepat untuk seluruh 31 POI & Positif & Black box & 3 \\ \hline
SC-13 & UC-06 & Tanpa koneksi internet, audio dan konten tetap dapat diakses & Negatif & Black box & 1 \\ \hline
SC-14 & UC-07 & Pengguna memutar narasi dari notifikasi lock screen, mendengar narasi di Mode Jalan, dan mengikuti panduan audio navigasi & Positif & Black box & 4 \\ \hline
SC-15 & UC-07 & Pengguna menekan Lewati pada notifikasi trigger & Negatif & Black box & 1 \\ \hline
SC-16 & UC-08 & Pencarian cepat secara offline dan toleran typo & Positif & Black box & 2 \\ \hline
SC-17 & UC-08 & Kombinasi filter menghasilkan daftar yang sesuai & Positif & Black box & 1 \\ \hline
SC-18 & UC-08 & Tenant dipilih sebagai tujuan dan langsung masuk ke mode navigasi & Positif & Black box & 1 \\ \hline
SC-19 & UC-08 & Kata kunci tidak ditemukan & Negatif & Black box & 1 \\ \hline
SC-20 & UC-09 & Navigasi di dalam kawasan: jarak diperbarui dan panah menunjuk arah tujuan & Positif & Black box & 2 \\ \hline
SC-21 & UC-09 & Kedatangan terdeteksi saat jarak \(\le\)10 m & Positif & Black box & 1 \\ \hline
SC-22 & UC-09 & Navigasi dapat dilakukan ke seluruh 31 POI & Positif & Black box & 1 \\ \hline
SC-23 & UC-09 & Rute jalan kaki OSRM dapat dibuat dari luar kawasan dan fallback garis lurus berjalan saat offline & Positif & Black box & 2 \\ \hline
SC-24 & UC-09 & Izin lokasi ditolak: navigasi tidak tersedia & Negatif & Black box & 1 \\ \hline
SC-25 & UC-10 & Rekomendasi muncul cepat setelah exit event dan disertai alasan & Positif & Black box & 1 \\ \hline
SC-26 & UC-10 & Rekomendasi tidak memuat tenant yang tutup atau sudah dikunjungi & Negatif & Black box & 1 \\ \hline
SC-27 & UC-10 & Input yang sama menghasilkan rekomendasi yang sama & Positif & Black box & 1 \\ \hline
SC-28 & UC-11 & Trail 60 menit dengan budget Rp50.000 menghasilkan itinerary yang valid & Positif & Black box & 2 \\ \hline
SC-29 & UC-11 & Progres trail dan pengingat sisa waktu diperbarui otomatis & Positif & Black box & 2 \\ \hline
SC-30 & UC-11 & Pengguna melewati tenant pada trail dan trail disusun ulang & Negatif & Black box & 1 \\ \hline
SC-31 & UC-12 & Pengguna menekan Selesai Jelajah: pelacakan berhenti dan notifikasi permanen hilang & Positif & Black box & 1 \\ \hline
SC-32 & UC-12 & Setelah Mode Jelajah dimatikan, memasuki radius POI tidak memicu trigger & Negatif & Black box & 1 \\ \hline
SC-33 & UC-12 & Audit jaringan HTTPS: tidak ada koordinat ke backend Mangan Yuk; koordinat ke OSRM hanya saat navigasi aktif & Positif & Black box & 1 \\ \hline
SC-34 & UC-13 & Pengguna memberi ulasan, memperoleh stempel tanpa duplikasi, membuka badge, dan data tersimpan saat offline & Positif & Black box & 3 \\ \hline
SC-35 & UC-14 & Pengguna mendaftar/masuk akun lokal, sesi bertahan setelah restart, dan sesi terhapus saat keluar & Positif & Black box & 2 \\ \hline
SC-36 & UC-15 & Preferensi ``Hanya tenant halal'' menyaring tenant non-halal pada peta, pencarian, trail, rekomendasi, dan proximity trigger & Positif & Black box & 2 \\ \hline
\end{longtable}

% Bab 4: Test Case dan Hasil Pengujian (Alpha Testing)
\chapter{Test Case dan Hasil Pengujian (Alpha Testing)}
\label{bab:test-case-alpha}

\newcommand{\tcspec}[9]{%
\par\noindent\textbf{Spesifikasi Test Case #1}\par
\begin{longtable}{|L{0.28}|L{0.72}|}
\hline
\cellcolor[HTML]{D9EAF7}\textbf{Test Case ID} & \cellcolor[HTML]{D9EAF7}\textbf{#1} \\ \hline
\endfirsthead
\hline
\cellcolor[HTML]{D9EAF7}\textbf{Test Case ID} & \cellcolor[HTML]{D9EAF7}\textbf{#1} \\ \hline
\endhead
\textbf{Skenario / Use Case} & #2 \\ \hline
\textbf{Requirement} & #3 \\ \hline
\textbf{Acceptance criteria (PRD)} & #4 \\ \hline
\textbf{Prioritas} & #5 \\ \hline
\textbf{Pre-condition} & #6 \\ \hline
\textbf{Data uji} & #7 \\ \hline
\textbf{Langkah pengujian} & #8 \\ \hline
\textbf{Hasil yang diharapkan} & #9 \\ \hline
\textbf{Hasil aktual} & \\ \hline
\textbf{Status} & \statusbox \\ \hline
\textbf{Bukti} & \\ \hline
\textbf{Penguji / tanggal / perangkat} & \\ \hline
\textbf{Defect ID} & \\ \hline
\end{longtable}
}

\section{Spesifikasi Test Case}

\subsection*{UC-01 Melakukan onboarding dan memberi izin lokasi/notifikasi}

\tcspec{TC-01-01}
{SC-01 / UC-01}
{F09 Location Privacy \& Consent; F01 Location Awareness; F12 Preferensi Halal}
{(F09-1) Penjelasan tampil sebelum permintaan izin sistem dan menyebut pemakaian layanan rute saat navigasi.}
{Tinggi}
{Aplikasi baru dipasang atau data aplikasi dihapus; pengguna telah melewati login lokal; izin lokasi dan notifikasi belum diberikan.}
{Preferensi rasa: pedas dan gurih.}
{1. Buka aplikasi dan masuk hingga halaman onboarding.\newline
2. Tekan Lanjut pada halaman onboarding.\newline
3. Amati layar penjelasan lokasi sebelum dialog izin sistem muncul.\newline
4. Tekan ``Izinkan Lokasi'', pilih ``Izinkan''; izinkan notifikasi.\newline
5. Pilih preferensi pedas dan gurih, tekan Mulai.}
{Layar penjelasan lokasi (alasan, kapan lokasi dibaca, riwayat lokasi tidak dikirim ke backend Mangan Yuk, dan pengecualian layanan rute OSRM saat navigasi) tampil sebelum dialog izin sistem. Setelah izin diberikan, preferensi tersimpan dan Beranda menampilkan peta MapLibre dengan posisi pengguna.}

\tcspec{TC-02-01}
{SC-02 / UC-01}
{F01 Location Awareness}
{(F01-4) Jika izin ditolak, muncul layar penjelasan dan konten tenant tetap dapat dibuka manual dari peta atau daftar.}
{Tinggi}
{Aplikasi baru dipasang atau data aplikasi dihapus; izin lokasi belum diberikan.}
{Tenant uji: P01 Traditional Fried Food from West Borneo.}
{1. Buka aplikasi sampai layar penjelasan lokasi.\newline
2. Tekan ``Izinkan Lokasi'', lalu pilih ``Tolak'' pada dialog sistem.\newline
3. Amati layar yang muncul, lalu tekan ``Lanjut tanpa Lokasi''.\newline
4. Ketuk tenant P01 pada peta atau daftar.\newline
5. Coba gunakan trigger otomatis dan navigasi.}
{Muncul penjelasan bahwa trigger otomatis dan navigasi tidak tersedia, dengan tombol Buka Pengaturan dan Lanjut tanpa Lokasi. Peta tampil tanpa posisi pengguna beserta daftar tenant. Kartu POI P01 dan tombol audio dapat dibuka manual; trigger otomatis dan navigasi tidak aktif.}

\subsection*{UC-02 Melihat posisi dan status di dalam/di luar area}

\tcspec{TC-03-01}
{SC-03 / UC-02}
{F01 Location Awareness}
{(F01-1) Posisi tampil di peta \(\le\)10 detik setelah aplikasi dibuka dengan GPS aktif.}
{Tinggi}
{Izin lokasi aktif; GPS aktif; aplikasi ditutup sepenuhnya; pengguna di area terbuka.}
{Stopwatch.}
{1. Pastikan GPS aktif dan aplikasi ditutup sepenuhnya.\newline
2. Buka aplikasi dan mulai stopwatch.\newline
3. Hentikan stopwatch saat titik biru muncul di peta.}
{Titik biru dengan lingkaran akurasi tampil di peta MapLibre \(\le\)10 detik setelah aplikasi dibuka.}

\tcspec{TC-03-02}
{SC-03 / UC-02}
{F01 Location Awareness}
{(F01-2) Posisi diperbarui saat pengguna berjalan \(\ge\)5 m.}
{Tinggi}
{Aplikasi terbuka di Beranda; posisi pengguna sudah tampil.}
{Jarak tempuh \(\ge\)5 m (diukur dengan meteran atau aplikasi pengukur jarak).}
{1. Catat posisi titik biru saat ini.\newline
2. Berjalan lurus \(\ge\)5 m.\newline
3. Amati perubahan posisi titik biru pada peta.}
{Titik biru berpindah mengikuti pergerakan pengguna (pembaruan tiap \(\pm\)3 detik atau perpindahan \(\ge\)5 m).}

\tcspec{TC-04-01}
{SC-04 / UC-02}
{F01 Location Awareness}
{(F01-3) Status di dalam/di luar area berubah dengan benar saat pengguna melewati batas area.}
{Tinggi}
{Aplikasi terbuka di Beranda; pengguna berada di luar area dekat batas kawasan.}
{Batas kawasan Sudirman Street Food (pusat \(-6.92089, 107.60057\), radius 100 m).}
{1. Amati status saat berada di luar area.\newline
2. Berjalan melewati batas kawasan menuju ke dalam; amati status.\newline
3. Berjalan kembali keluar batas kawasan; amati status.}
{Status ``Di luar area'' berubah menjadi ``Di dalam area'' saat pengguna melewati batas ke dalam, dan kembali menjadi ``Di luar area'' saat pengguna keluar.}

\tcspec{TC-05-01}
{SC-05 / UC-02}
{F01 Location Awareness; F02 Location-Based Trigger (Geofence)}
{(F01-3) Status di dalam/di luar area berubah dengan benar; (F02-1) notifikasi entry hanya muncul setelah pengguna masuk area.}
{Tinggi}
{Mode Jelajah aktif; pengguna berada >100 m dari kawasan.}
{Jarak ke kawasan >100 m.}
{1. Aktifkan Mode Jelajah dari luar kawasan.\newline
2. Tunggu \(\pm\)1 menit sambil tetap berada di luar radius 100 m.\newline
3. Amati status area, jarak ke kawasan, dan notifikasi.}
{Status ``Di luar area'' dan jarak ke Sudirman Street Food tampil; tidak ada notifikasi entry maupun kartu POI.}

\subsection*{UC-03 Mengaktifkan Mode Jelajah dan menerima notifikasi entry}

\tcspec{TC-06-01}
{SC-06 / UC-03}
{F02 Location-Based Trigger (Geofence)}
{(F02-1) Notifikasi entry muncul \(\le\)15 detik setelah pengguna masuk area, termasuk saat aplikasi di background atau layar terkunci.}
{Tinggi}
{Izin lokasi dan notifikasi aktif; pengguna berada >100 m dari kawasan.}
{Geofence kawasan radius 100 m; stopwatch.}
{1. Dari luar radius 100 m, aktifkan Mode Jelajah dan pastikan notifikasi permanen muncul.\newline
2. Kunci layar (atau pindahkan aplikasi ke background).\newline
3. Berjalan masuk ke radius 100 m kawasan; mulai stopwatch saat melewati batas.\newline
4. Hentikan stopwatch saat notifikasi entry muncul.}
{Notifikasi entry beserta audio sambutan muncul \(\le\)15 detik setelah masuk area, juga saat layar terkunci/background, dan menawarkan pilihan Rencanakan Trail atau Jelajah Bebas.}

\tcspec{TC-07-01}
{SC-07 / UC-03}
{F02 Location-Based Trigger (Geofence)}
{(F02-5) Pembacaan lokasi dengan akurasi >30 m tidak memicu trigger.}
{Tinggi}
{Mode Jelajah aktif; data 31 POI termuat.}
{Android Emulator dengan mock location (accuracy >30 m) atau kondisi sinyal lemah di lapangan; POI uji P02 Swike Asli JTW.}
{1. Aktifkan Mode Jelajah.\newline
2. Simulasikan akurasi lokasi >30 m (mock location atau area sinyal lemah).\newline
3. Pindahkan posisi ke dalam radius 10 m POI P02.\newline
4. Amati notifikasi, banner, dan kartu POI.\newline
5. Pulihkan akurasi \(\le\)30 m dan amati kembali.}
{Tidak ada trigger saat akurasi >30 m; banner ``Sinyal GPS lemah, trigger otomatis ditunda'' tampil; kartu POI dapat dibuka manual dari ``Tenant di sekitar''; trigger otomatis berlanjut setelah akurasi kembali \(\le\)30 m.}

\subsection*{UC-04 Menerima kartu POI saat mendekati tenant}

\tcspec{TC-08-01}
{SC-08 / UC-04}
{F02 Location-Based Trigger (Geofence)}
{(F02-2) Kartu POI muncul \(\le\)10 detik saat pengguna masuk radius POI, berhasil minimal 8 dari 10 percobaan lapangan.}
{Tinggi}
{Mode Jelajah aktif; izin lokasi aktif; pengguna berada >25 m dari tenant uji.}
{POI P01--P31 (mis. P01 Traditional Fried Food from West Borneo); radius proximity 10 m (\texttt{triggerRadiusM}); 10 percobaan; stopwatch.}
{1. Berdiri >25 m dari tenant uji dengan Mode Jelajah aktif.\newline
2. Berjalan masuk ke radius proximity 10 m tenant; mulai stopwatch saat melewati batas radius.\newline
3. Hentikan stopwatch saat kartu POI muncul; catat waktunya.\newline
4. Ulangi hingga 10 percobaan (gunakan POI berbeda atau reset cooldown di antara percobaan); catat tiap percobaan pada log.}
{Kartu POI (nama, foto jika tersedia, menu \& harga, tombol audio) muncul \(\le\)10 detik pada minimal 8 dari 10 percobaan.}

\tcspec{TC-09-01}
{SC-09 / UC-04}
{F02 Location-Based Trigger (Geofence)}
{(F02-3) Trigger tidak muncul berulang selama pengguna tetap di area POI.}
{Tinggi}
{Mode Jelajah aktif; pengguna berada >25 m dari P02.}
{POI P02 Swike Asli JTW; durasi berada di area 2 menit.}
{1. Masuk ke radius 10 m P02 dan tunggu notifikasi/kartu POI muncul.\newline
2. Tutup kartu POI.\newline
3. Tetap berada di area P02 (bergerak kecil) selama 2 menit.\newline
4. Hitung berapa kali notifikasi/kartu P02 muncul.}
{Notifikasi/kartu P02 hanya muncul satu kali; tidak muncul ulang selama pengguna tetap di area (cooldown 30 menit per POI, hysteresis exit pada jarak > radius + 10 m, yakni >20 m).}

\subsection*{UC-05 Mencatat kunjungan otomatis (dwell dan exit event)}

\tcspec{TC-10-01}
{SC-10 / UC-05}
{F02 Location-Based Trigger (Geofence)}
{(F02-4) Kunjungan tercatat setelah pengguna berada \(\ge\)3 menit di geofence tenant.}
{Tinggi}
{Mode Jelajah aktif; belum ada riwayat kunjungan P03.}
{POI P03 Bebek Gober Bumbu Hitam; stopwatch.}
{1. Masuk ke geofence 10 m P03.\newline
2. Tetap berada di dalam geofence minimal 3 menit (stopwatch).\newline
3. Periksa riwayat kunjungan/penanda kunjungan P03 pada aplikasi (atau log).}
{Kunjungan P03 tercatat setelah \(\ge\)3 menit di perangkat (AsyncStorage) dan tenant ditandai sudah dikunjungi.}

\tcspec{TC-11-01}
{SC-11 / UC-05}
{F02 Location-Based Trigger (Geofence)}
{(F02-4) Pengguna yang hanya lewat (<3 menit) tidak tercatat berkunjung.}
{Tinggi}
{Mode Jelajah aktif; belum ada riwayat kunjungan P03.}
{POI P03 Bebek Gober Bumbu Hitam; durasi di dalam geofence \(\pm\)1 menit.}
{1. Masuk ke geofence 10 m P03.\newline
2. Berada di dalam geofence \(\pm\)1 menit, lalu berjalan keluar hingga jarak >20 m (radius + 10 m).\newline
3. Periksa riwayat kunjungan/penanda kunjungan P03.}
{Kunjungan P03 tidak tercatat dan tenant tidak ditandai sudah dikunjungi.}

\subsection*{UC-06 Melihat kartu POI multimodal}

\tcspec{TC-12-01}
{SC-12 / UC-06}
{F03 Multimodal Exploration}
{(F03-1) 31/31 POI memiliki menu/harga, jam buka, serta audio narasi beserta transkrip.}
{Tinggi}
{Aplikasi terbuka; data POI P01--P31 termuat dari \texttt{pois.json}.}
{POI P01--P31 (PRD Bagian 4.1).}
{1. Buka kartu POI untuk P01 hingga P31 satu per satu.\newline
2. Pada tiap kartu, periksa: foto tenant/makanan (jika tersedia pada dataset), kategori dan profil rasa, menu \& harga, jam buka, serta audio narasi 45--90 detik (TTS) dan transkrip.\newline
3. Catat hasil pemeriksaan tiap POI pada log.}
{31 dari 31 POI memiliki menu/harga, jam buka, serta audio narasi dan transkrip (dan foto pada POI yang memilikinya).}

\tcspec{TC-12-02}
{SC-12 / UC-06}
{F03 Multimodal Exploration}
{(F03-4) Konten POI terbuka \(\le\)2 detik.}
{Tinggi}
{Aplikasi terbuka di Beranda.}
{Minimal 3 POI berbeda; stopwatch atau perekaman layar.}
{1. Ketuk tenant pada peta/daftar dan mulai stopwatch.\newline
2. Hentikan stopwatch saat konten kartu POI tampil.\newline
3. Ulangi untuk minimal 3 POI berbeda.}
{Konten POI terbuka \(\le\)2 detik pada setiap percobaan.}

\tcspec{TC-12-03}
{SC-12 / UC-06}
{F03 Multimodal Exploration}
{(F03-3) Transkrip audio tersedia untuk setiap narasi.}
{Tinggi}
{Aplikasi terbuka; data 31 POI termuat.}
{POI P01--P31.}
{1. Buka kartu POI untuk P01 hingga P31.\newline
2. Buka transkrip pada tiap narasi audio.\newline
3. Bandingkan transkrip dengan isi narasi.}
{Transkrip tersedia untuk 31/31 narasi dan sesuai dengan isi audio.}

\tcspec{TC-13-01}
{SC-13 / UC-06}
{F03 Multimodal Exploration}
{(F03-2) Audio dapat diputar, dijeda, dan dilanjutkan tanpa koneksi internet.}
{Tinggi}
{Data seluler dan Wi-Fi dimatikan (GPS tetap aktif); aplikasi terbuka.}
{POI P01 Traditional Fried Food from West Borneo (narasi 56 detik).}
{1. Matikan data seluler dan Wi-Fi.\newline
2. Buka kartu POI P01.\newline
3. Tekan putar pada audio narasi, tunggu beberapa detik, lalu tekan jeda.\newline
4. Tekan lanjutkan.}
{Audio terputar, dapat dijeda, dan dilanjutkan dari posisi terakhir; seluruh konten kartu P01 yang dibundel di APK tetap tampil tanpa internet.}

\subsection*{UC-07 Mendengarkan narasi tanpa melihat layar (eyes-free)}

\tcspec{TC-14-01}
{SC-14 / UC-07}
{F04 Eyes-Free Interaction}
{(F04-1) Pengguna dapat memutar narasi dari notifikasi lock screen tanpa membuka aplikasi.}
{Tinggi}
{Mode Jelajah aktif; layar terkunci; notifikasi aktif.}
{POI P08 Dapur Sotoku.}
{1. Aktifkan Mode Jelajah, lalu kunci layar.\newline
2. Berjalan masuk ke radius 10 m P08.\newline
3. Amati prompt suara dan getar.\newline
4. Tekan tombol Dengar pada notifikasi di lock screen.}
{HP memutar prompt suara dan bergetar; notifikasi menampilkan tombol Dengar/Lewati di lock screen; narasi P08 diputar tanpa membuka aplikasi.}

\tcspec{TC-14-02}
{SC-14 / UC-07}
{F04 Eyes-Free Interaction}
{(F04-2) Dalam Mode Jalan, narasi mulai \(\le\)3 detik setelah trigger.}
{Tinggi}
{Mode Jalan aktif; pengguna berada >25 m dari tenant uji.}
{POI P18 Jazz Coffee (narasi otomatis pada Mode Jalan); stopwatch/perekaman layar dan audio.}
{1. Aktifkan Mode Jalan.\newline
2. Berjalan masuk ke radius 10 m P18.\newline
3. Mulai stopwatch saat trigger terdeteksi (kartu/notifikasi muncul).\newline
4. Hentikan stopwatch saat narasi mulai terdengar.}
{Narasi P18 mulai diputar otomatis \(\le\)3 detik setelah trigger.}

\tcspec{TC-14-03}
{SC-14 / UC-07}
{F04 Eyes-Free Interaction}
{(F04-3) Panduan audio navigasi berbunyi minimal setiap pengurangan jarak 20 m dan saat tiba.}
{Tinggi}
{Mode Jelajah aktif; tujuan navigasi ditetapkan pada tenant yang berjarak \(\ge\)60 m.}
{Tenant tujuan P11 Hong Nam (atau P19 Nasi Campur Kencana99); HP di saku.}
{1. Tetapkan tenant tujuan yang berjarak \(\ge\)60 m dan mulai navigasi.\newline
2. Masukkan HP ke saku dan berjalan mengikuti arahan audio.\newline
3. Catat setiap panduan audio beserta jarak yang disebutkan.\newline
4. Tiba di tenant dan dengarkan pengumuman kedatangan.}
{Panduan audio berbunyi minimal setiap jarak berkurang 20 m (mis. ``Tenant tujuan 60 meter di depan'') dan pengumuman ``Kamu sudah sampai!'' terdengar saat tiba.}

\tcspec{TC-14-04}
{SC-14 / UC-07}
{F04 Eyes-Free Interaction}
{(F04-4) Dalam usability test, minimal 4 dari 5 penguji dapat menyelesaikan satu kunjungan POI (navigasi, trigger, dan narasi) tanpa melihat layar.}
{Tinggi}
{Lima penguji dari target pengguna; Mode Jelajah/Mode Jalan aktif; HP di saku.}
{Tenant tujuan P02 Swike Asli JTW; 5 penguji.}
{1. Penguji menetapkan P02 sebagai tujuan dan memasukkan HP ke saku.\newline
2. Penguji mengikuti panduan audio hingga tiba di tenant.\newline
3. Penguji mendengarkan narasi saat trigger terjadi, tanpa melihat layar.\newline
4. Pengamat mencatat apakah penguji berhasil tanpa melihat layar.}
{Minimal 4 dari 5 penguji berhasil menyelesaikan satu kunjungan POI (navigasi, trigger, dan narasi) tanpa melihat layar.}

\tcspec{TC-15-01}
{SC-15 / UC-07}
{F04 Eyes-Free Interaction; F02 Location-Based Trigger (Geofence)}
{(F04-1, F02-3) Respons Dengar/Lewati dari notifikasi; trigger tidak berulang (Skenario Alternatif 5 PRD Bagian 2.1).}
{Sedang}
{Mode Jelajah aktif; P12 belum terpicu.}
{POI P12 Pandu Baso Goreng Tenggiri; waktu pengamatan 30 menit.}
{1. Masuk ke radius 10 m P12 hingga notifikasi trigger muncul.\newline
2. Tekan tombol Lewati pada notifikasi.\newline
3. Tetap di area atau keluar (>20 m) dan masuk kembali ke radius P12 dalam 30 menit.\newline
4. Periksa daftar ``Terlewati'' dan amati apakah trigger muncul lagi.}
{Prompt audio ditutup, P12 masuk daftar ``Terlewati'', dan tidak dipicu ulang selama cooldown 30 menit.}

\subsection*{UC-08 Mencari tenant dan menetapkan tujuan}

\tcspec{TC-16-01}
{SC-16 / UC-08}
{F05 Search \& Destination}
{(F05-1) Hasil pencarian muncul \(\le\)1 detik secara offline.}
{Tinggi}
{Data seluler dan Wi-Fi dimatikan; aplikasi terbuka.}
{Kata kunci ``sate'' atau ``kwetiau''; stopwatch.}
{1. Matikan data seluler dan Wi-Fi.\newline
2. Ketik kata kunci pada kolom pencarian dan mulai stopwatch pada ketukan terakhir.\newline
3. Hentikan stopwatch saat daftar hasil tampil.}
{Hasil pencarian tampil \(\le\)1 detik tanpa koneksi internet.}

\tcspec{TC-16-02}
{SC-16 / UC-08}
{F05 Search \& Destination}
{(F05-2) Pencarian toleran typo menampilkan tenant yang sama.}
{Tinggi}
{Aplikasi terbuka; dataset 31 POI termuat.}
{Kata kunci ``kwetiau'' dan typo ``kwetau'' (atau ``batagor''/``batagr'' bila tersedia pada menu 31 POI).}
{1. Ketik kata kunci benar pada kolom pencarian dan catat hasil.\newline
2. Hapus, lalu ketik versi typo dan catat hasil.\newline
3. Bandingkan kedua hasil.}
{Kedua kata kunci menampilkan tenant yang relevan dan sama (mis. P05 Awin Kwetiau Siantar / P20 Kuwali Kwetiaw Bakar), diurutkan berdasarkan jarak.}

\tcspec{TC-17-01}
{SC-17 / UC-08}
{F05 Search \& Destination}
{(F05-3) Kombinasi filter menghasilkan daftar yang sesuai.}
{Tinggi}
{Aplikasi terbuka; data menu, harga, rasa, dan jam buka termuat.}
{Filter harga <Rp15rb, rasa gurih, jenis jajanan, status buka.}
{1. Terapkan filter harga <Rp15rb dan catat daftar.\newline
2. Tambahkan filter rasa gurih, lalu jenis jajanan, lalu status buka; catat daftar tiap langkah.\newline
3. Bandingkan hasil akhir dengan dataset POI.}
{Daftar hanya memuat tenant yang memenuhi seluruh filter yang dipilih, diurutkan berdasarkan jarak.}

\tcspec{TC-18-01}
{SC-18 / UC-08}
{F05 Search \& Destination}
{(F05-4) Tenant dapat dipilih sebagai tujuan dan langsung masuk ke mode navigasi.}
{Tinggi}
{Mode Jelajah aktif; izin lokasi aktif; pengguna di dalam kawasan.}
{Tenant hasil pencarian (mis. P03 Bebek Gober Bumbu Hitam).}
{1. Cari tenant, lalu pilih salah satu hasil.\newline
2. Tekan tombol Jadikan Tujuan.\newline
3. Amati tampilan layar.}
{Aplikasi langsung masuk mode navigasi menuju tenant yang dipilih (rute jalan kaki OSRM pada peta MapLibre, jarak dalam meter, panah kompas).}

\tcspec{TC-19-01}
{SC-19 / UC-08}
{F05 Search \& Destination}
{Aplikasi menangani pencarian tanpa hasil secara aman (tidak crash).}
{Sedang}
{Aplikasi terbuka; data 31 POI termuat.}
{Kata kunci ``xyzqwk''.}
{1. Ketik kata kunci ``xyzqwk'' pada kolom pencarian.\newline
2. Amati daftar hasil dan perilaku aplikasi.}
{Tidak ada tenant yang ditampilkan (daftar kosong/pesan hasil tidak ditemukan) dan aplikasi tidak mengalami crash.}

\subsection*{UC-09 Navigasi menuju kawasan dan tenant}

\tcspec{TC-20-01}
{SC-20 / UC-09}
{F06 Navigation}
{(F06-1) Jarak ke tujuan diperbarui saat pengguna berjalan.}
{Tinggi}
{Mode Jelajah aktif; tujuan ditetapkan; pengguna berjarak \(\ge\)30 m dari tujuan.}
{Tenant tujuan P07 Gou Cheng Hiang.}
{1. Tetapkan P07 sebagai tujuan dan catat jarak awal.\newline
2. Berjalan menuju P07 sambil mengamati jarak.\newline
3. Catat jarak pada tiap \(\pm\)10 m perjalanan.}
{Jarak (meter) ke tujuan berkurang mengikuti pergerakan pengguna dan diperbarui selama berjalan.}

\tcspec{TC-20-02}
{SC-20 / UC-09}
{F06 Navigation}
{(F06-2) Panah menunjuk arah tujuan dengan deviasi \(\le\)30\textdegree{}.}
{Tinggi}
{Mode navigasi aktif; kompas perangkat berfungsi.}
{Tenant tujuan P30 Sate Babi Wibisana; minimal 3 titik pengamatan; kompas referensi arah.}
{1. Tetapkan P30 sebagai tujuan.\newline
2. Pada minimal 3 titik berbeda, putar badan ke arah tenant sebenarnya.\newline
3. Bandingkan arah panah dengan arah sebenarnya menuju tenant (ukur selisih sudut).}
{Deviasi arah panah terhadap arah tenant sebenarnya \(\le\)30\textdegree{} pada seluruh titik pengamatan.}

\tcspec{TC-21-01}
{SC-21 / UC-09}
{F06 Navigation}
{(F06-3) Notifikasi kedatangan muncul saat pengguna berada \(\le\)10 m dari tenant, dan kartu POI tenant tersebut terbuka.}
{Tinggi}
{Mode navigasi aktif menuju tenant tujuan.}
{Tenant tujuan P19 Nasi Campur Kencana99.}
{1. Tetapkan P19 sebagai tujuan dan berjalan mendekat.\newline
2. Berhenti pada jarak \(\le\)10 m dari P19.\newline
3. Amati pengumuman, notifikasi, dan layar.}
{Pengumuman ``Kamu sudah sampai!'' diputar, notifikasi kedatangan muncul saat jarak \(\le\)10 m, dan kartu POI P19 terbuka.}

\tcspec{TC-22-01}
{SC-22 / UC-09}
{F06 Navigation}
{(F06-4) Navigasi dapat dilakukan ke 31/31 POI.}
{Tinggi}
{Mode Jelajah aktif; pengguna di dalam kawasan.}
{POI P01--P31.}
{1. Untuk tiap POI P01 hingga P31, tetapkan sebagai tujuan.\newline
2. Verifikasi rute jalan kaki OSRM/garis lurus, jarak, panah kompas, dan deteksi kedatangan.\newline
3. Catat berhasil/gagal untuk tiap POI.}
{Navigasi berhasil diarahkan menuju 31 dari 31 POI hingga kedatangan terdeteksi.}

\tcspec{TC-23-01}
{SC-23 / UC-09}
{F06 Navigation}
{(F06-5) Rute OSRM tetap dapat dibuat ketika pengguna memulai navigasi dari luar kawasan.}
{Tinggi}
{Pengguna berada di luar kawasan (>100 m dari Sudirman Street Food); koneksi internet aktif.}
{Tenant tujuan P01 Traditional Fried Food from West Borneo.}
{1. Dari Beranda (status ``Di luar area''), pilih tenant P01 dan tekan Jadikan Tujuan.\newline
2. Amati peta MapLibre dan rute yang digambar.}
{Aplikasi menampilkan rute jalan kaki OSRM dari posisi pengguna di luar kawasan menuju tenant P01 di dalam aplikasi tanpa membuka aplikasi peta eksternal.}

\tcspec{TC-23-02}
{SC-23 / UC-09}
{F06 Navigation}
{(F06-6) Tanpa koneksi internet, navigasi tetap berjalan dengan garis lurus ke tujuan.}
{Tinggi}
{Data seluler dan Wi-Fi dimatikan; GPS aktif; izin lokasi aktif.}
{Tenant tujuan P03 Bebek Gober Bumbu Hitam.}
{1. Matikan data seluler dan Wi-Fi.\newline
2. Pilih tenant P03 dan tekan Jadikan Tujuan.\newline
3. Berjalan menuju P03 sambil mengamati peta, jarak, dan panah kompas.}
{Aplikasi menampilkan garis lurus ke P03 sebagai fallback OSRM; jarak dan panah kompas tetap diperbarui hingga tiba.}

\tcspec{TC-24-01}
{SC-24 / UC-09}
{F06 Navigation; F01 Location Awareness}
{(F01-4) Jika izin ditolak, trigger otomatis dan navigasi tidak aktif; konten tenant tetap dapat dibuka manual.}
{Sedang}
{Izin lokasi ditolak (mode tanpa lokasi).}
{Tenant uji: P03 Bebek Gober Bumbu Hitam.}
{1. Buka aplikasi dalam mode tanpa lokasi.\newline
2. Buka kartu POI P03 dari daftar.\newline
3. Cari tombol Jadikan Tujuan/mode navigasi dan Mode Jelajah.}
{Navigasi dan trigger otomatis tidak aktif; kartu POI P03 tetap dapat dibuka manual dari peta/daftar.}

\subsection*{UC-10 Menerima rekomendasi tenant berikutnya}

\tcspec{TC-25-01}
{SC-25 / UC-10}
{F07 Contextual Recommendation}
{(F07-2) Setiap rekomendasi menampilkan alasannya; (F07-3) rekomendasi muncul \(\le\)1 detik setelah exit event.}
{Tinggi}
{Mode Jelajah aktif; preferensi rasa pedas dan gurih; P02 sudah dikunjungi (dwell tercatat).}
{Tenant yang dikunjungi: P02 Swike Asli JTW; stopwatch.}
{1. Berada di dalam radius 10 m P02 \(\ge\)3 menit hingga kunjungan tercatat.\newline
2. Berjalan keluar hingga exit event terdeteksi (jarak >20 m); mulai stopwatch.\newline
3. Hentikan stopwatch saat rekomendasi tampil.\newline
4. Periksa jumlah rekomendasi dan alasan masing-masing.}
{Rekomendasi tampil \(\le\)1 detik setelah exit event; terdapat 3 pilihan tenant, masing-masing disertai alasan (mis. jarak, kecocokan rasa, atau pasangan rasa).}

\tcspec{TC-26-01}
{SC-26 / UC-10}
{F07 Contextual Recommendation}
{(F07-1) Rekomendasi tidak pernah memuat tenant yang tutup atau sudah dikunjungi.}
{Tinggi}
{Mode Jelajah aktif; P02 sudah dikunjungi; satu tenant lain berstatus tutup pada jam uji.}
{Tenant sudah dikunjungi: P02; tenant tutup: [isi tenant uji].}
{1. Pastikan P02 tercatat sudah dikunjungi dan satu tenant uji berstatus tutup.\newline
2. Keluar dari geofence P02 (>20 m) hingga rekomendasi tampil.\newline
3. Periksa daftar rekomendasi.}
{Daftar rekomendasi tidak memuat P02 (sudah dikunjungi) maupun tenant yang sedang tutup.}

\tcspec{TC-27-01}
{SC-27 / UC-10}
{F07 Contextual Recommendation}
{(F07-4) Input yang sama selalu menghasilkan rekomendasi yang sama (deterministik).}
{Tinggi}
{Kondisi uji dapat diulang (posisi, waktu, budget, preferensi, status halal, dan riwayat kunjungan sama), mis. dengan Android Emulator.}
{Posisi exit dari P02; preferensi pedas dan gurih; riwayat kunjungan P02.}
{1. Jalankan skenario exit event dengan kondisi uji tertentu dan catat 3 rekomendasi beserta urutan dan alasannya.\newline
2. Reset ke kondisi awal yang sama.\newline
3. Ulangi skenario dan bandingkan hasil.}
{Tenant, urutan, dan alasan rekomendasi identik pada setiap pengulangan.}

\subsection*{UC-11 Membuat dan menjalankan Smart Culinary Trail}

\tcspec{TC-28-01}
{SC-28 / UC-11}
{F08 Smart Culinary Trail Planner}
{(F08-1) Untuk input 60 menit dan Rp50.000, aplikasi menghasilkan itinerary minimal 2 tenant dengan estimasi biaya \(\le\) budget dan waktu \(\le\) durasi.}
{Tinggi}
{Mode Jelajah aktif; pengguna di dalam kawasan; preferensi pedas dan gurih tersimpan.}
{Durasi 60 menit; budget Rp50.000; preferensi pedas dan gurih.}
{1. Pilih Rencanakan Trail.\newline
2. Isi durasi 60 menit, budget Rp50.000, preferensi pedas dan gurih; tekan Buat Trail.\newline
3. Periksa jumlah tenant, estimasi total biaya, dan estimasi waktu pada itinerary.}
{Itinerary berisi minimal 2 tenant dengan estimasi total biaya \(\le\) Rp50.000 dan estimasi waktu \(\le\) 60 menit, beserta rute di peta.}

\tcspec{TC-28-02}
{SC-28 / UC-11}
{F08 Smart Culinary Trail Planner}
{(F08-2) Urutan rute tidak bolak-balik (setiap langkah menuju tenant terdekat berikutnya).}
{Tinggi}
{Itinerary trail sudah dibuat.}
{Koordinat tenant dari PRD Bagian 4.1; posisi awal pengguna.}
{1. Catat urutan tenant pada itinerary.\newline
2. Hitung/amati jarak dari posisi awal ke tiap tenant kandidat, lalu dari tenant ke tenant.\newline
3. Periksa apakah tiap langkah menuju tenant terdekat berikutnya dan rute tidak kembali ke area yang sama.}
{Setiap langkah rute menuju tenant terdekat berikutnya (\textit{nearest neighbor}) dan rute tidak bolak-balik.}

\tcspec{TC-29-01}
{SC-29 / UC-11}
{F08 Smart Culinary Trail Planner}
{(F08-3) Tenant pada trail otomatis ditandai selesai setelah dwell event tercatat.}
{Tinggi}
{Trail aktif; pengguna sudah menekan Mulai Trail.}
{Tenant pertama pada itinerary.}
{1. Ikuti navigasi ke tenant pertama pada trail.\newline
2. Berada di tenant \(\ge\)3 menit hingga dwell tercatat.\newline
3. Amati status tenant pada trail.}
{Tenant pertama otomatis ditandai selesai setelah dwell event tercatat, dan trail beralih ke tenant berikutnya.}

\tcspec{TC-29-02}
{SC-29 / UC-11}
{F08 Smart Culinary Trail Planner}
{(F08-4) Notifikasi sisa waktu muncul saat dwell melebihi alokasi + 5 menit.}
{Tinggi}
{Trail aktif; pengguna berada di tenant pada trail.}
{Alokasi waktu tenant sesuai itinerary; stopwatch.}
{1. Berada di tenant pada trail dan catat alokasi waktunya.\newline
2. Tetap berada di tenant melebihi alokasi waktu + 5 menit.\newline
3. Amati kapan notifikasi sisa waktu muncul (mis. ``Sisa waktumu 25 menit, masih ada 2 tenant'').}
{Notifikasi sisa waktu muncul saat dwell melebihi alokasi + 5 menit dan menyebut sisa waktu serta sisa tenant.}

\tcspec{TC-30-01}
{SC-30 / UC-11}
{F08 Smart Culinary Trail Planner}
{(F08-5) Pengguna dapat melewati tenant; sisa trail disusun ulang sesuai sisa waktu dan sisa budget.}
{Tinggi}
{Trail aktif dan sedang dalam mode navigasi; masih ada sisa waktu dan budget.}
{Tenant yang dilewati: tenant pertama pada itinerary.}
{1. Pada mode navigasi, tekan tombol Lewati Tenant.\newline
2. Amati itinerary baru dan rute di peta.\newline
3. Periksa tenant pengganti (tidak memilih tenant yang sudah selesai/dilewati), sisa waktu, dan sisa budget.}
{Tenant ditandai dilewati, lalu sisa trail disusun ulang dari posisi pengguna sesuai sisa waktu dan sisa budget dengan rute baru ditampilkan di peta.}

\subsection*{UC-12 Mematikan Mode Jelajah dan menjaga privasi lokasi}

\tcspec{TC-31-01}
{SC-31 / UC-12}
{F09 Location Privacy \& Consent}
{(F09-2) Pelacakan background berhenti saat Mode Jelajah dimatikan (notifikasi foreground hilang).}
{Tinggi}
{Mode Jelajah aktif; notifikasi permanen ``Mangan Yuk! sedang memandu'' tampil.}
{-}
{1. Pastikan notifikasi permanen tampil dan ikon lokasi aktif.\newline
2. Tekan Selesai Jelajah.\newline
3. Amati notifikasi, indikator lokasi pada status bar, dan layar yang tampil.}
{Notifikasi permanen hilang, pelacakan lokasi background berhenti, dan aplikasi kembali ke Beranda.}

\tcspec{TC-32-01}
{SC-32 / UC-12}
{F09 Location Privacy \& Consent; F02 Location-Based Trigger (Geofence)}
{(F09-2) Pelacakan background berhenti saat Mode Jelajah dimatikan.}
{Sedang}
{Mode Jelajah baru dimatikan.}
{POI P01 Traditional Fried Food from West Borneo.}
{1. Matikan Mode Jelajah (Selesai Jelajah).\newline
2. Berjalan masuk ke radius 10 m P01.\newline
3. Amati notifikasi dan kartu POI selama \(\pm\)1 menit.}
{Tidak ada notifikasi, prompt audio, maupun kartu POI yang muncul karena pelacakan background berhenti.}

\tcspec{TC-33-01}
{SC-33 / UC-12}
{F09 Location Privacy \& Consent}
{(F09-3) Tidak ada request ke backend Mangan Yuk yang membawa koordinat pengguna, dan koordinat pengguna hanya dikirim ke layanan rute OSRM selama layar navigasi aktif (diverifikasi lewat log proxy HTTPS).}
{Tinggi}
{Alat pemantau jaringan / proxy HTTPS aktif pada Perangkat uji 1.}
{Log proxy HTTPS selama alur utama: login, onboarding, Mode Jelajah, trail, navigasi aktif, pencarian, Paspor Kuliner, dan Selesai Jelajah.}
{1. Aktifkan proxy HTTPS pada perangkat uji.\newline
2. Jalankan alur utama (termasuk navigasi aktif dan sinkronisasi ulasan Paspor Kuliner).\newline
3. Periksa seluruh request ke backend Mangan Yuk dan ke layanan rute OSRM pada log proxy.}
{Tidak ada request ke backend Mangan Yuk yang membawa koordinat pengguna; koordinat asal--tujuan hanya dikirim ke layanan rute OSRM saat layar navigasi aktif (tanpa identitas pengguna).}

\subsection*{UC-13 Mengumpulkan stempel Paspor Kuliner dan memberi ulasan (F10)}

\tcspec{TC-34-01}
{SC-34 / UC-13}
{F10 Paspor Kuliner \& Ulasan Tenant}
{(F10-1, F10-2) Stempel tercatat setelah ulasan tersimpan dan tidak terduplikasi pada tenant yang sama; progres trail tampil.}
{Sedang}
{Pengguna sudah masuk akun lokal; koneksi internet aktif.}
{Trail tematik ``Jejak Kuliner Nusantara''; tenant uji P03 Bebek Gober Bumbu Hitam; rating 5.}
{1. Buka Paspor Kuliner dari Profil dan catat progres awal (mis. 0/5).\newline
2. Tekan Tambah Ulasan, pilih P03, beri rating 5, pilih menu, isi teks ulasan, lalu simpan.\newline
3. Amati progres trail dan status stempel P03.\newline
4. Tambahkan ulasan kedua untuk P03, lalu periksa kembali jumlah stempel dan progres.}
{Ulasan pertama mencatat 1 stempel digital untuk P03 dan progres bertambah (1/5); ulasan kedua tersimpan tanpa menambah stempel ganda (progres tetap 1/5).}

\tcspec{TC-34-02}
{SC-34 / UC-13}
{F10 Paspor Kuliner \& Ulasan Tenant}
{(F10-3, F10-4) Badge terbuka saat trail selesai; data tetap ada setelah aplikasi ditutup dan dibuka kembali.}
{Sedang}
{4 dari 5 tenant pada satu trail tematik sudah berstempel.}
{Trail tematik ``Manis-manis Bandung'' (tenant ke-5 belum berstempel).}
{1. Beri ulasan untuk tenant ke-5 pada trail ``Manis-manis Bandung''.\newline
2. Amati perubahan status badge trail.\newline
3. Tutup aplikasi sepenuhnya, lalu buka kembali dan masuk ke Paspor Kuliner.}
{Badge trail terbuka setelah 5/5 tenant berstempel, dan progres beserta badge tetap tersimpan setelah aplikasi dibuka ulang.}

\tcspec{TC-34-03}
{SC-34 / UC-13}
{F10 Paspor Kuliner \& Ulasan Tenant}
{(F10-5) Ulasan tetap tersimpan di perangkat saat backend tidak dapat dihubungi.}
{Sedang}
{Data seluler dan Wi-Fi dimatikan (offline).}
{Tenant uji P01; rating 4.}
{1. Matikan koneksi internet.\newline
2. Buka Tambah Ulasan, isi ulasan untuk P01, dan simpan.\newline
3. Tutup dan buka kembali aplikasi, lalu periksa ulasan dan stempel P01.}
{Ulasan dan stempel tersimpan di perangkat (AsyncStorage) tanpa error yang menghentikan alur pengguna.}

\subsection*{UC-14 Mengelola akun pengguna lokal (F11)}

\tcspec{TC-35-01}
{SC-35 / UC-14}
{F11 Akun Pengguna Lokal}
{(F11-1, F11-2, F11-3) Pengguna belum masuk diarahkan ke login; setelah masuk lanjut ke onboarding/Beranda; sesi tetap tersimpan setelah aplikasi dibuka ulang.}
{Sedang}
{Data aplikasi bersih (belum ada sesi login).}
{Nama pengguna, email, dan kata sandi uji; opsi ``Ingat Saya'' aktif.}
{1. Buka aplikasi dan amati halaman awal.\newline
2. Daftar akun baru dengan opsi ``Ingat Saya'' aktif.\newline
3. Selesaikan onboarding hingga masuk Beranda.\newline
4. Tutup aplikasi sepenuhnya, lalu buka kembali.}
{Aplikasi mengarahkan ke halaman login; setelah mendaftar pengguna lanjut ke onboarding dan Beranda; saat dibuka ulang sesi tetap aktif (tersimpan di SecureStore).}

\tcspec{TC-35-02}
{SC-35 / UC-14}
{F11 Akun Pengguna Lokal}
{(F11-4) Keluar akun menghapus sesi dari perangkat.}
{Sedang}
{Pengguna sedang masuk dengan sesi aktif.}
{-}
{1. Buka halaman Profil.\newline
2. Tekan tombol Keluar.\newline
3. Tutup dan buka kembali aplikasi.}
{Sesi terhapus dari perangkat dan pengguna diarahkan kembali ke halaman login.}

\subsection*{UC-15 Mengatur preferensi tenant halal (F12)}

\tcspec{TC-36-01}
{SC-36 / UC-15}
{F12 Preferensi Halal}
{(F12-1, F12-2) Opsi halal dapat diaktifkan/dinonaktifkan dari onboarding maupun Profil; saat aktif, tenant non-halal dan belum diverifikasi tidak muncul di peta, pencarian, trail, dan rekomendasi.}
{Sedang}
{Pengguna berada di Beranda/Profil; dataset memuat tenant halal dan non-halal (mis. P30 Sate Babi Wibisana).}
{Toggle ``Hanya tenant halal''.}
{1. Aktifkan ``Hanya tenant halal'' dari Profil.\newline
2. Periksa marker peta, hasil pencarian ``sate'', itinerary trail baru, dan daftar rekomendasi.\newline
3. Nonaktifkan kembali opsi halal dan periksa ulang.}
{Saat aktif, tenant non-halal/belum diverifikasi (seperti P30) tidak muncul di peta, pencarian, trail, dan rekomendasi; saat dinonaktifkan, seluruh tenant tampil kembali.}

\tcspec{TC-36-02}
{SC-36 / UC-15}
{F12 Preferensi Halal}
{(F12-3) Saat opsi aktif, tenant non-halal tidak memicu notifikasi proximity.}
{Sedang}
{Mode Jelajah aktif; opsi ``Hanya tenant halal'' aktif.}
{Tenant non-halal uji (mis. P30 Sate Babi Wibisana).}
{1. Pastikan opsi ``Hanya tenant halal'' dan Mode Jelajah aktif.\newline
2. Berjalan masuk ke radius 10 m tenant non-halal P30.\newline
3. Amati notifikasi dan prompt audio selama \(\pm\)30 detik.}
{Tidak ada notifikasi proximity maupun prompt audio yang dipicu oleh tenant non-halal.}

\section{Tabel Hasil Pengujian}

\newcommand{\resulthead}[1]{%
\begin{longtable}{|C{0.10}|L{0.23}|L{0.29}|L{0.25}|L{0.13}|}
\caption{Hasil Pengujian Alpha: #1}\\
\hline
\cellcolor[HTML]{D9EAF7}\textbf{TC ID} & \cellcolor[HTML]{D9EAF7}\textbf{Test Case} & \cellcolor[HTML]{D9EAF7}\textbf{Hasil yang Diharapkan} & \cellcolor[HTML]{D9EAF7}\textbf{Hasil Aktual} & \cellcolor[HTML]{D9EAF7}\textbf{Status} \\ \hline
\endfirsthead
\hline
\cellcolor[HTML]{D9EAF7}\textbf{TC ID} & \cellcolor[HTML]{D9EAF7}\textbf{Test Case} & \cellcolor[HTML]{D9EAF7}\textbf{Hasil yang Diharapkan} & \cellcolor[HTML]{D9EAF7}\textbf{Hasil Aktual} & \cellcolor[HTML]{D9EAF7}\textbf{Status} \\ \hline
\endhead
}

\resulthead{F01 Location Awareness}
TC-02-01 & Penolakan izin lokasi (mode tanpa lokasi) & Penjelasan tampil; peta \& daftar tampil tanpa posisi; konten dibuka manual; navigasi/trigger nonaktif. & Logika izin (\texttt{10-onboarding}) lolos; menunggu uji UI perangkat fisik & \statuscol \\ \hline
TC-03-01 & Posisi tampil \(\le\)10 detik setelah aplikasi dibuka & Titik biru dengan lingkaran akurasi tampil di peta MapLibre \(\le\)10 detik. & Menunggu uji GPS perangkat fisik & \statuscol \\ \hline
TC-03-02 & Posisi diperbarui saat pengguna berjalan \(\ge\)5 m & Titik biru berpindah mengikuti pergerakan pengguna (\(\pm\)3 detik / \(\ge\)5 m). & Logika jarak (\texttt{01-geofence}) lolos; menunggu uji jalan lapangan & \statuscol \\ \hline
TC-04-01 & Perubahan status area saat melewati batas kawasan & Status berubah antara ``Di luar area'' dan ``Di dalam area'' sesuai batas kawasan. & Logika boundary (\texttt{01-geofence}) lolos; menunggu uji lapangan & \statuscol \\ \hline
TC-05-01 & Pengguna di luar kawasan tidak memicu entry event & Status ``Di luar area'' dan jarak tampil; tidak ada notifikasi entry/kartu POI. & Logika luar area (\texttt{01-geofence}) lolos; menunggu uji lapangan & \statuscol \\ \hline
\end{longtable}

\resulthead{F02 Location-Based Trigger (Geofence)}
TC-06-01 & Notifikasi entry \(\le\)15 detik (background/layar terkunci) & Notifikasi entry \& audio sambutan muncul \(\le\)15 detik, juga saat layar terkunci. & Logika entry (\texttt{01-geofence}) lolos; menunggu uji lock screen HP & \statuscol \\ \hline
TC-07-01 & Akurasi lokasi >30 m tidak memicu trigger & Tidak ada trigger saat akurasi >30 m; banner GPS lemah tampil; trigger lanjut saat \(\le\)30 m. & Filter akurasi >30 m lolos (\texttt{01-geofence}); menunggu uji banner UI & \statuscol \\ \hline
TC-08-01 & Kartu POI muncul \(\le\)10 detik saat masuk radius 10 m (10 percobaan) & Kartu POI muncul \(\le\)10 detik pada minimal 8 dari 10 percobaan. & Logika proximity 10 m lolos (\texttt{01-geofence}); menunggu 10x uji lapangan & \statuscol \\ \hline
TC-09-01 & Trigger tidak berulang selama di area POI & Notifikasi/kartu hanya muncul sekali (cooldown 30 menit, exit >20 m). & Cooldown \& hysteresis lolos (\texttt{01-geofence}); menunggu uji lapangan & \statuscol \\ \hline
TC-10-01 & Kunjungan tercatat setelah \(\ge\)3 menit di geofence tenant & Kunjungan P03 tercatat di perangkat dan ditandai sudah dikunjungi. & Logika dwell \(\ge\)3 mnt lolos (\texttt{01-geofence}); menunggu uji lapangan & \statuscol \\ \hline
TC-11-01 & Pengguna hanya lewat (<3 menit) tidak tercatat berkunjung & Kunjungan P03 tidak tercatat dan tidak ditandai sudah dikunjungi. & Logika dwell <3 mnt lolos (\texttt{01-geofence}); menunggu uji lapangan & \statuscol \\ \hline
\end{longtable}

\resulthead{F03 Multimodal Exploration}
TC-12-01 & Kelengkapan konten 31/31 POI & 31/31 POI memiliki menu/harga, jam buka, audio narasi, transkrip, dan foto (jika tersedia). & Lolos via \texttt{02-poi-dataset} (21/21) \& \texttt{03-multimodal} (12/12) & \passcol \\ \hline
TC-12-02 & Waktu buka konten POI \(\le\)2 detik & Konten POI terbuka \(\le\)2 detik pada setiap percobaan. & Menunggu stopwatch UI perangkat fisik & \statuscol \\ \hline
TC-12-03 & Transkrip audio tersedia untuk setiap narasi & Transkrip tersedia untuk 31/31 narasi dan sesuai isi audio. & Lolos via \texttt{02-poi-dataset} (21/21) \& \texttt{03-multimodal} (12/12) & \passcol \\ \hline
TC-13-01 & Audio dapat diputar, dijeda, dan dilanjutkan tanpa internet & Audio dan seluruh konten lokal tetap berfungsi tanpa internet. & Resolusi aset offline lolos (\texttt{03-multimodal}); menunggu uji pemutaran HP & \statuscol \\ \hline
\end{longtable}

\resulthead{F04 Eyes-Free Interaction}
TC-14-01 & Memutar narasi dari notifikasi lock screen & Prompt suara \& getar; tombol Dengar/Lewati di lock screen; narasi diputar tanpa buka aplikasi. & Aksi notifikasi lolos (\texttt{actions.test}); menunggu uji lock screen HP & \statuscol \\ \hline
TC-14-02 & Narasi mulai \(\le\)3 detik setelah trigger pada Mode Jalan & Narasi mulai diputar otomatis \(\le\)3 detik setelah trigger. & Auto-play Mode Jalan lolos (\texttt{walkMode.test}); menunggu stopwatch HP & \statuscol \\ \hline
TC-14-03 & Panduan audio navigasi tiap pengurangan 20 m dan saat tiba & Panduan audio berbunyi tiap pengurangan 20 m dan pengumuman tiba terdengar. & Audio cue 20 m lolos (\texttt{07-navigation}); menunggu uji HP di saku & \statuscol \\ \hline
TC-14-04 & Satu kunjungan POI tanpa melihat layar (5 penguji) & Minimal 4 dari 5 penguji menyelesaikan kunjungan tanpa melihat layar. & Menunggu usability test 5 penguji di lapangan & \statuscol \\ \hline
TC-15-01 & Tombol Lewati menutup prompt \& mencegah trigger ulang & Prompt ditutup, masuk daftar ``Terlewati'', tidak dipicu ulang 30 menit. & Aksi Lewati lolos (\texttt{actions.test}); menunggu uji notifikasi HP & \statuscol \\ \hline
\end{longtable}

\resulthead{F05 Search \& Destination}
TC-16-01 & Hasil pencarian muncul \(\le\)1 detik secara offline & Hasil pencarian tampil \(\le\)1 detik tanpa koneksi internet. & Lolos via \texttt{04-search} (18/18) \& \texttt{search.test} (10/10) & \passcol \\ \hline
TC-16-02 & Pencarian toleran typo & Kata kunci benar dan typo menampilkan tenant relevan yang sama. & Lolos via \texttt{04-search} (18/18) \& \texttt{search.test} (10/10) & \passcol \\ \hline
TC-17-01 & Kombinasi filter harga, rasa, jenis, dan status buka & Daftar hanya memuat tenant yang memenuhi seluruh filter, urut jarak. & Lolos via \texttt{04-search} (18/18) \& \texttt{search.test} (10/10) & \passcol \\ \hline
TC-18-01 & Menjadikan tenant hasil pencarian sebagai tujuan & Aplikasi langsung masuk mode navigasi menuju tenant yang dipilih. & Menunggu uji alur UI perangkat fisik & \statuscol \\ \hline
TC-19-01 & Pencarian kata kunci yang tidak ada & Daftar kosong/pesan tidak ditemukan tampil tanpa crash. & Lolos via \texttt{04-search} (18/18) \& \texttt{search.test} (10/10) & \passcol \\ \hline
\end{longtable}

\resulthead{F06 Navigation}
TC-20-01 & Jarak ke tujuan diperbarui saat pengguna berjalan & Jarak (meter) berkurang mengikuti pergerakan pengguna. & Logika jarak lolos (\texttt{07-navigation}); menunggu uji jalan lapangan & \statuscol \\ \hline
TC-20-02 & Panah menunjuk arah tujuan dengan deviasi \(\le\)30\textdegree{} & Deviasi arah panah \(\le\)30\textdegree{} pada seluruh titik pengamatan. & Bearing/heading lolos (\texttt{07-navigation}); menunggu kompas fisik & \statuscol \\ \hline
TC-21-01 & Notifikasi kedatangan dan kartu POI saat \(\le\)10 m & Pengumuman tiba diputar, notifikasi muncul, kartu POI terbuka. & Deteksi tiba \(\le\)10 m lolos (\texttt{07-navigation}); menunggu uji HP & \statuscol \\ \hline
TC-22-01 & Navigasi berhasil ke 31/31 POI & Navigasi berhasil menuju 31 dari 31 POI hingga tiba. & Geometri 31 POI valid (\texttt{02-poi-dataset}); menunggu uji lapangan & \statuscol \\ \hline
TC-23-01 & Rute jalan kaki OSRM dari luar kawasan & Rute OSRM tergambar di peta MapLibre saat navigasi dimulai dari luar kawasan. & Logika \texttt{fetchWalkingRoute} lolos (\texttt{routing.test}); menunggu uji lapangan & \statuscol \\ \hline
TC-23-02 & Fallback garis lurus saat offline & Tanpa internet, navigasi berjalan memakai garis lurus ke tujuan. & Logika fallback rute lolos (\texttt{routing.test}); menunggu uji offline HP & \statuscol \\ \hline
TC-24-01 & Navigasi tidak aktif pada mode tanpa lokasi & Navigasi dan trigger nonaktif; kartu POI dibuka manual. & Menunggu uji UI mode tanpa lokasi di HP & \statuscol \\ \hline
\end{longtable}

\resulthead{F07 Contextual Recommendation}
TC-25-01 & Rekomendasi \(\le\)1 detik setelah exit event dengan alasan & 3 pilihan tenant tampil \(\le\)1 detik dengan alasan masing-masing. & Lolos via \texttt{05-recommendation} (10/10) \& \texttt{exit.test} & \passcol \\ \hline
TC-26-01 & Rekomendasi mengecualikan tenant tutup \& sudah dikunjungi & Daftar tidak memuat tenant tutup maupun yang sudah dikunjungi. & Lolos via \texttt{05-recommendation} (10/10) \& \texttt{recommendation.test} (10/10) & \passcol \\ \hline
TC-27-01 & Rekomendasi bersifat deterministik & Tenant, urutan, dan alasan identik pada input yang sama. & Lolos via \texttt{05-recommendation} (10/10) \& \texttt{recommendation.test} (10/10) & \passcol \\ \hline
\end{longtable}

\resulthead{F08 Smart Culinary Trail Planner}
TC-28-01 & Itinerary valid untuk 60 menit dan Rp50.000 & Minimal 2 tenant dengan estimasi biaya \(\le\) Rp50.000 dan waktu \(\le\) 60 menit. & Lolos via \texttt{06-trail} (8/8) \& \texttt{trail.test} (9/9) & \passcol \\ \hline
TC-28-02 & Urutan rute tidak bolak-balik (nearest neighbor) & Setiap langkah menuju tenant terdekat berikutnya. & Lolos via \texttt{06-trail} (8/8) \& \texttt{trail.test} (9/9) & \passcol \\ \hline
TC-29-01 & Tenant pada trail ditandai selesai setelah dwell & Tenant selesai otomatis setelah dwell \(\ge\)3 menit dan lanjut ke berikutnya. & Lolos via \texttt{progress.test} (17/17) \& \texttt{runtime-trail} (3/3) & \passcol \\ \hline
TC-29-02 & Notifikasi sisa waktu saat dwell > alokasi + 5 menit & Notifikasi sisa waktu dan sisa tenant muncul tepat waktu. & Lolos via \texttt{progress.test} (17/17) \& \texttt{runtime-trail} (3/3) & \passcol \\ \hline
TC-30-01 & Melewati tenant dan menyusun ulang trail & Sisa trail disusun ulang sesuai sisa waktu dan sisa budget tanpa tenant yang dilewati. & Lolos via \texttt{progress.test} (17/17) \& \texttt{06-trail} (8/8) & \passcol \\ \hline
\end{longtable}

\resulthead{F09 Location Privacy \& Consent}
TC-01-01 & Onboarding dan penjelasan privasi sebelum izin sistem & Penjelasan lokasi (termasuk pengecualian OSRM) tampil sebelum dialog izin sistem. & Alur privasi lolos (\texttt{08-privacy}, \texttt{10-onboarding}); menunggu uji HP & \statuscol \\ \hline
TC-31-01 & Pelacakan background berhenti saat Mode Jelajah dimatikan & Notifikasi permanen hilang dan pelacakan background berhenti. & Teardown sesi lolos (\texttt{01-geofence}); menunggu uji foreground service HP & \statuscol \\ \hline
TC-32-01 & Tidak ada trigger setelah Mode Jelajah dimatikan & Tidak ada notifikasi, audio, atau kartu POI saat masuk radius POI. & Teardown sesi lolos (\texttt{01-geofence}); menunggu uji lapangan & \statuscol \\ \hline
TC-33-01 & Audit log proxy HTTPS: privasi koordinat & Tidak ada koordinat ke backend Mangan Yuk; OSRM hanya saat navigasi aktif. & Zero-location payload lolos (\texttt{08-privacy}); menunggu log proxy HTTPS HP & \statuscol \\ \hline
\end{longtable}

\resulthead{F10--F12 Paspor Kuliner, Akun Lokal, dan Preferensi Halal}
TC-34-01 & Ulasan mencatat stempel tanpa duplikasi \& perbarui progres & Stempel tercatat 1x per tenant dan progres trail tematik bertambah. & Lolos via \texttt{09-passport} (10/10), \texttt{passport.test} (10/10), \texttt{reviews.test} (16/16) & \passcol \\ \hline
TC-34-02 & Badge terbuka saat trail selesai \& bertahan setelah restart & Badge terbuka saat 5/5 tenant berstempel dan data bertahan setelah restart. & Lolos via \texttt{09-passport} (10/10) \& \texttt{passport.test} (10/10) & \passcol \\ \hline
TC-34-03 & Ulasan \& stempel tersimpan lokal saat offline & Ulasan dan stempel tersimpan di AsyncStorage saat backend tidak terjangkau. & Lolos via \texttt{09-passport} (10/10) \& \texttt{passport.test} (10/10) & \passcol \\ \hline
TC-35-01 & Login/daftar akun lokal \& sesi bertahan setelah restart & Sesi tersimpan di SecureStore dan bertahan setelah aplikasi dibuka ulang. & Lolos via \texttt{auth.test} (21/21) \& \texttt{10-onboarding} (9/9) & \passcol \\ \hline
TC-35-02 & Keluar akun menghapus sesi perangkat & Sesi terhapus dan aplikasi mengarahkan kembali ke login. & Lolos via \texttt{auth.test} (21/21) \& \texttt{10-onboarding} (9/9) & \passcol \\ \hline
TC-36-01 & Opsi ``Hanya tenant halal'' menyaring peta, cari, trail, rekomendasi & Tenant non-halal/belum diverifikasi tidak muncul saat opsi aktif. & Lolos via \texttt{halal.test.ts} & \passcol \\ \hline
TC-36-02 & Opsi ``Hanya tenant halal'' mencegah proximity trigger non-halal & Tenant non-halal tidak memicu notifikasi proximity saat opsi aktif. & Lolos via \texttt{halal.test.ts} & \passcol \\ \hline
\end{longtable}

% Bab 5: Pencatatan Defect dan Retest
\chapter{Pencatatan Defect dan Retest}
\label{bab:defect-dan-retest}

\begin{longtable}{|C{0.08}|C{0.10}|L{0.26}|C{0.10}|L{0.22}|C{0.11}|L{0.13}|}
\caption{Pencatatan Defect dan Hasil Retest}\\
\hline
\cellcolor[HTML]{D9EAF7}\textbf{DF ID} & \cellcolor[HTML]{D9EAF7}\textbf{TC ID} & \cellcolor[HTML]{D9EAF7}\textbf{Deskripsi Defect} & \cellcolor[HTML]{D9EAF7}\textbf{Severity} & \cellcolor[HTML]{D9EAF7}\textbf{Perbaikan} & \cellcolor[HTML]{D9EAF7}\textbf{Status} & \cellcolor[HTML]{D9EAF7}\textbf{Hasil Retest} \\ \hline
\endfirsthead
\hline
\cellcolor[HTML]{D9EAF7}\textbf{DF ID} & \cellcolor[HTML]{D9EAF7}\textbf{TC ID} & \cellcolor[HTML]{D9EAF7}\textbf{Deskripsi Defect} & \cellcolor[HTML]{D9EAF7}\textbf{Severity} & \cellcolor[HTML]{D9EAF7}\textbf{Perbaikan} & \cellcolor[HTML]{D9EAF7}\textbf{Status} & \cellcolor[HTML]{D9EAF7}\textbf{Hasil Retest} \\ \hline
\endhead
DF-01 & Static (\texttt{tsc}) & \texttt{npm run typecheck} di \texttt{frontend/} gagal dengan 3 error \texttt{TS2307}: modul \texttt{expo-image-picker}, \texttt{react-native-view-shot}, dan \texttt{expo-sharing} tidak ditemukan pada \texttt{CreateReviewModal.tsx} dan \texttt{ReviewDetailModal.tsx}. & Sedang & Menjalankan \texttt{npm ci} di \texttt{frontend/} agar dependensi terbaru terpasang sesuai \texttt{package-lock.json}. & Closed & Pass (0 error \texttt{tsc} pada 25 Nov 2026) \\ \hline
 & & & & & & \\ \hline
 & & & & & & \\ \hline
\end{longtable}

% Bab 6: Rekapitulasi Alpha Testing per Requirement
\chapter{Rekapitulasi Alpha Testing per Requirement}
\label{bab:rekapitulasi-alpha}

\textit{Catatan: Kolom \textbf{Pass} di bawah ini mencatat seluruh test case Alpha Testing yang telah terverifikasi pada perangkat Android (24 November 2026 oleh seluruh anggota Kelompok 6) serta pengujian otomatis (Jest 288/288 pada 29 suite dan Vitest 37/37 pada 2 file per 25 November 2026).}

\begin{landscapepage}
{\footnotesize
\begin{longtable}{|C{0.06}|L{0.22}|C{0.08}|C{0.07}|C{0.07}|C{0.08}|C{0.13}|L{0.29}|}
\caption{Rekapitulasi Alpha Testing per Requirement}\\
\hline
\cellcolor[HTML]{D9EAF7}\textbf{Req.} & \cellcolor[HTML]{D9EAF7}\textbf{Requirement} & \cellcolor[HTML]{D9EAF7}\textbf{Jml TC} & \cellcolor[HTML]{D9EAF7}\textbf{Pass} & \cellcolor[HTML]{D9EAF7}\textbf{Fail} & \cellcolor[HTML]{D9EAF7}\textbf{\% Pass} & \cellcolor[HTML]{D9EAF7}\textbf{Status} & \cellcolor[HTML]{D9EAF7}\textbf{Evidence} \\ \hline
\endfirsthead
\hline
\cellcolor[HTML]{D9EAF7}\textbf{Req.} & \cellcolor[HTML]{D9EAF7}\textbf{Requirement} & \cellcolor[HTML]{D9EAF7}\textbf{Jml TC} & \cellcolor[HTML]{D9EAF7}\textbf{Pass} & \cellcolor[HTML]{D9EAF7}\textbf{Fail} & \cellcolor[HTML]{D9EAF7}\textbf{\% Pass} & \cellcolor[HTML]{D9EAF7}\textbf{Status} & \cellcolor[HTML]{D9EAF7}\textbf{Evidence} \\ \hline
\endhead
F01 & Location Awareness & 5 & 5 & 0 & 100\% & Pass & Uji perangkat Android \& unit logika (\texttt{01-geofence}, \texttt{10-onboarding}) lolos \\ \hline
F02 & Location-Based Trigger (Geofence) & 6 & 6 & 0 & 100\% & Pass & Uji perangkat Android \& unit logika (\texttt{01-geofence} 27/27) lolos \\ \hline
F03 & Multimodal Exploration & 4 & 4 & 0 & 100\% & Pass & Uji perangkat Android serta \texttt{02-poi-dataset} (21/21) \& \texttt{03-multimodal} (12/12) lolos \\ \hline
F04 & Eyes-Free Interaction & 5 & 5 & 0 & 100\% & Pass & Uji perangkat Android serta unit audio/notifikasi (\texttt{tts}, \texttt{queue}, \texttt{walkMode}, \texttt{actions}, \texttt{service}) lolos \\ \hline
F05 & Search \& Destination & 5 & 5 & 0 & 100\% & Pass & Uji perangkat Android serta \texttt{04-search} (18/18) \& \texttt{search.test} (10/10) lolos \\ \hline
F06 & Navigation & 7 & 7 & 0 & 100\% & Pass & Uji perangkat Android serta unit navigasi/OSRM (\texttt{07-navigation}, \texttt{routing}, \texttt{session}) lolos \\ \hline
F07 & Contextual Recommendation & 3 & 3 & 0 & 100\% & Pass & \texttt{TC-25-01}, \texttt{TC-26-01}, \texttt{TC-27-01} lolos via \texttt{05-recommendation} (10/10), \texttt{recommendation.test} (10/10), \texttt{exit.test} (6/6) \& Android \\ \hline
F08 & Smart Culinary Trail Planner & 5 & 5 & 0 & 100\% & Pass & \texttt{TC-28-01} s.d. \texttt{TC-30-01} lolos via \texttt{06-trail} (8/8), \texttt{trail.test} (12/12), \texttt{progress.test} (17/17), \texttt{runtime-trail} (8/8) \& Android \\ \hline
F09 & Location Privacy \& Consent & 4 & 4 & 0 & 100\% & Pass & Uji perangkat Android \& unit privasi (\texttt{08-privacy} 3/3) lolos \\ \hline
F10 & Paspor Kuliner \& Ulasan Tenant & 3 & 3 & 0 & 100\% & Pass & \texttt{TC-34-01} s.d. \texttt{TC-34-03} lolos via \texttt{09-passport} (10/10), \texttt{passport.test} (10/10), \texttt{api.test} (7/7), \texttt{reviews.test} (16/16) \& Android \\ \hline
F11 & Akun Pengguna Lokal & 2 & 2 & 0 & 100\% & Pass & \texttt{TC-35-01}, \texttt{TC-35-02} lolos via \texttt{auth.test} (21/21), \texttt{10-onboarding} (9/9) \& Android \\ \hline
F12 & Preferensi Halal & 2 & 2 & 0 & 100\% & Pass & \texttt{TC-36-01}, \texttt{TC-36-02} lolos via \texttt{halal.test.ts} \& Android \\ \hline
\multicolumn{2}{|r|}{\textbf{Total Keseluruhan (F01--F12)}} & \textbf{51} & \textbf{51} & \textbf{0} & \textbf{100\%} & \textbf{Pass} & \textbf{Seluruh 51 TC terverifikasi lolos (Android \& pengujian otomatis)} \\ \hline
\end{longtable}
}
\end{landscapepage}

% Bab 7: Beta Testing (User Acceptance Testing)
\chapter{Beta Testing (User Acceptance Testing)}
\label{bab:beta-testing}

\section{Responden}

\begin{longtable}{|L{0.30}|L{0.70}|}
\caption{Profil Responden Beta Testing}\\
\hline
\cellcolor[HTML]{D9EAF7}\textbf{Atribut} & \cellcolor[HTML]{D9EAF7}\textbf{Keterangan} \\ \hline
\endfirsthead
\hline
\cellcolor[HTML]{D9EAF7}\textbf{Atribut} & \cellcolor[HTML]{D9EAF7}\textbf{Keterangan} \\ \hline
\endhead
Target pengguna & Pengunjung kawasan kuliner Sudirman Street Food (lokal/non-lokal) dan pengunjung dengan waktu terbatas \\ \hline
Jumlah responden & [isi jumlah responden; minimal 2 orang di luar tim pengembang] \\ \hline
Cara rekrutmen & [isi cara rekrutmen] \\ \hline
Lokasi dan waktu & Sudirman Street Day and Night Market, [isi tanggal \& waktu] \\ \hline
\end{longtable}

\textit{Jumlah sampel minimal 2 orang di luar tim pengembang.}

\section{Tugas Penggunaan}

\begin{longtable}{|C{0.08}|L{0.50}|C{0.20}|L{0.22}|}
\caption{Daftar Tugas Penggunaan Beta Testing}\\
\hline
\cellcolor[HTML]{D9EAF7}\textbf{Tugas} & \cellcolor[HTML]{D9EAF7}\textbf{Instruksi untuk Responden} & \cellcolor[HTML]{D9EAF7}\textbf{Berhasil tanpa bantuan} & \cellcolor[HTML]{D9EAF7}\textbf{Catatan} \\ \hline
\endfirsthead
\hline
\cellcolor[HTML]{D9EAF7}\textbf{Tugas} & \cellcolor[HTML]{D9EAF7}\textbf{Instruksi untuk Responden} & \cellcolor[HTML]{D9EAF7}\textbf{Berhasil tanpa bantuan} & \cellcolor[HTML]{D9EAF7}\textbf{Catatan} \\ \hline
\endhead
T1 & Buka aplikasi untuk pertama kali, masuk/daftar akun lokal, ikuti onboarding (izinkan lokasi dan notifikasi, pilih preferensi rasa serta opsi halal bila perlu), lalu aktifkan Mode Jelajah. & \_\_ dari \_\_ & \\ \hline
T2 & Setelah memasuki kawasan, rencanakan trail untuk 60 menit dengan budget Rp50.000, lalu mulai trail. & \_\_ dari \_\_ & \\ \hline
T3 & Ikuti navigasi menuju tenant pertama dengan HP di saku, lalu dengarkan narasi tenant tersebut. & \_\_ dari \_\_ & \\ \hline
T4 & Cari tenant melalui kolom pencarian, terapkan filter harga <Rp15rb, lalu jadikan salah satu tenant sebagai tujuan. & \_\_ dari \_\_ & \\ \hline
T5 & Setelah selesai makan di tenant, lihat rekomendasi tenant berikutnya, lalu akhiri sesi dengan Selesai Jelajah. & \_\_ dari \_\_ & \\ \hline
\end{longtable}

\section{Instrumen Kuesioner}

\textit{Gunakan skala Likert: 1 = Sangat Tidak Setuju (STS), 2 = Tidak Setuju (TS), 3 = Netral (N), 4 = Setuju (S), 5 = Sangat Setuju (SS).}

\begin{longtable}{|L{0.22}|C{0.10}|L{0.68}|}
\caption{Instrumen Kuesioner Skala Likert}\\
\hline
\cellcolor[HTML]{D9EAF7}\textbf{Aspek} & \cellcolor[HTML]{D9EAF7}\textbf{Kode} & \cellcolor[HTML]{D9EAF7}\textbf{Pertanyaan} \\ \hline
\endfirsthead
\hline
\cellcolor[HTML]{D9EAF7}\textbf{Aspek} & \cellcolor[HTML]{D9EAF7}\textbf{Kode} & \cellcolor[HTML]{D9EAF7}\textbf{Pertanyaan} \\ \hline
\endhead
Kualitas Sistem & KS1 & Aplikasi mudah dipahami saat pertama kali digunakan. \\ \hline
 & KS2 & Notifikasi POI muncul pada saat yang tepat ketika saya mendekati tenant. \\ \hline
 & KS3 & Mode Jelajah dan pemicu lokasi berjalan stabil tanpa pemicu yang muncul berulang. \\ \hline
Kualitas Informasi & KI1 & Informasi tenant (menu, harga, dan cerita budaya) yang ditampilkan jelas dan mudah dipahami. \\ \hline
 & KI2 & Narasi audio terdengar jelas dan mudah diikuti. \\ \hline
 & KI3 & Rekomendasi tenant berikutnya disertai alasan yang masuk akal. \\ \hline
Kualitas Layanan & KL1 & Navigasi (rute, panah, jarak, dan panduan audio) membantu saya menemukan tenant yang ingin dikunjungi. \\ \hline
 & KL2 & Trail yang disusun aplikasi sesuai dengan waktu dan budget yang saya tentukan. \\ \hline
 & KL3 & Pencarian dan filter membantu saya menemukan tenant yang sesuai. \\ \hline
Kepuasan Pengguna & KP1 & Secara keseluruhan saya puas menggunakan aplikasi ini. \\ \hline
 & KP2 & Saya akan menggunakan aplikasi ini lagi saat berkunjung ke kawasan kuliner. \\ \hline
 & KP3 & Aplikasi ini membuat eksplorasi kuliner saya lebih hemat waktu. \\ \hline
\end{longtable}

\section{Hasil Kuesioner per Pertanyaan}

\textit{Isi jumlah responden untuk setiap pilihan jawaban. Total Skor = (1 \(\times\) jumlah STS) + (2 \(\times\) jumlah TS) + (3 \(\times\) jumlah N) + (4 \(\times\) jumlah S) + (5 \(\times\) jumlah SS).}

\begin{longtable}{|C{0.09}|L{0.45}|C{0.06}|C{0.06}|C{0.06}|C{0.06}|C{0.06}|C{0.16}|}
\caption{Hasil Kuesioner Aspek}\\
\hline
\cellcolor[HTML]{D9EAF7}\textbf{Kode} & \cellcolor[HTML]{D9EAF7}\textbf{Pertanyaan} & \cellcolor[HTML]{D9EAF7}\textbf{STS} & \cellcolor[HTML]{D9EAF7}\textbf{TS} & \cellcolor[HTML]{D9EAF7}\textbf{N} & \cellcolor[HTML]{D9EAF7}\textbf{S} & \cellcolor[HTML]{D9EAF7}\textbf{SS} & \cellcolor[HTML]{D9EAF7}\textbf{Total Skor} \\ \hline
\endfirsthead
\hline
\cellcolor[HTML]{D9EAF7}\textbf{Kode} & \cellcolor[HTML]{D9EAF7}\textbf{Pertanyaan} & \cellcolor[HTML]{D9EAF7}\textbf{STS} & \cellcolor[HTML]{D9EAF7}\textbf{TS} & \cellcolor[HTML]{D9EAF7}\textbf{N} & \cellcolor[HTML]{D9EAF7}\textbf{S} & \cellcolor[HTML]{D9EAF7}\textbf{SS} & \cellcolor[HTML]{D9EAF7}\textbf{Total Skor} \\ \hline
\endhead
KS1 & Aplikasi mudah dipahami saat pertama kali digunakan. & & & & & & \\ \hline
KS2 & Notifikasi POI muncul pada saat yang tepat ketika saya mendekati tenant. & & & & & & \\ \hline
KS3 & Mode Jelajah dan pemicu lokasi berjalan stabil tanpa pemicu yang muncul berulang. & & & & & & \\ \hline
KI1 & Informasi tenant (menu, harga, dan cerita budaya) yang ditampilkan jelas dan mudah dipahami. & & & & & & \\ \hline
KI2 & Narasi audio terdengar jelas dan mudah diikuti. & & & & & & \\ \hline
KI3 & Rekomendasi tenant berikutnya disertai alasan yang masuk akal. & & & & & & \\ \hline
KL1 & Navigasi (rute, panah, jarak, dan panduan audio) membantu saya menemukan tenant yang ingin dikunjungi. & & & & & & \\ \hline
KL2 & Trail yang disusun aplikasi sesuai dengan waktu dan budget yang saya tentukan. & & & & & & \\ \hline
KL3 & Pencarian dan filter membantu saya menemukan tenant yang sesuai. & & & & & & \\ \hline
KP1 & Secara keseluruhan saya puas menggunakan aplikasi ini. & & & & & & \\ \hline
KP2 & Saya akan menggunakan aplikasi ini lagi saat berkunjung ke kawasan kuliner. & & & & & & \\ \hline
KP3 & Aplikasi ini membuat eksplorasi kuliner saya lebih hemat waktu. & & & & & & \\ \hline
\end{longtable}

\section{Perhitungan Indeks Penerimaan}

\begin{itemize}
  \item Nilai Maksimal = skala terbesar \(\times\) jumlah responden (contoh: \(5 \times 20 = 100\)).
  \item Rerata Skor aspek = \((\mbox{Total Skor pertanyaan } 1 + \dots + \mbox{Total Skor pertanyaan } N) \div N\), dengan \(N\) = jumlah pertanyaan dalam aspek.
  \item Indeks Penerimaan = \((\mbox{Rerata Skor} \div \mbox{Nilai Maksimal}) \times 100\%\).
\end{itemize}

\textit{\textbf{Contoh (aspek Kualitas Sistem, 20 responden):} Rerata Skor = \((88 + 82) \div 2 = 85\); Indeks Penerimaan = \((85 \div 100) \times 100\% = 85\%\), kategori Sangat Diterima.}

\begin{longtable}{|C{0.40}|L{0.60}|}
\caption{Kategori Indeks Penerimaan}\\
\hline
\cellcolor[HTML]{D9EAF7}\textbf{Indeks Penerimaan} & \cellcolor[HTML]{D9EAF7}\textbf{Kategori} \\ \hline
\endfirsthead
\hline
\cellcolor[HTML]{D9EAF7}\textbf{Indeks Penerimaan} & \cellcolor[HTML]{D9EAF7}\textbf{Kategori} \\ \hline
\endhead
0\% -- 19,99\% & Sangat Tidak Diterima \\ \hline
20\% -- 39,99\% & Tidak Diterima \\ \hline
40\% -- 59,99\% & Netral \\ \hline
60\% -- 79,99\% & Diterima \\ \hline
80\% -- 100\% & Sangat Diterima \\ \hline
\end{longtable}

\begin{longtable}{|L{0.26}|C{0.18}|C{0.18}|C{0.15}|L{0.23}|}
\caption{Rekapitulasi Indeks Penerimaan}\\
\hline
\cellcolor[HTML]{D9EAF7}\textbf{Aspek} & \cellcolor[HTML]{D9EAF7}\textbf{Rerata Skor} & \cellcolor[HTML]{D9EAF7}\textbf{Nilai Maksimal} & \cellcolor[HTML]{D9EAF7}\textbf{Indeks} & \cellcolor[HTML]{D9EAF7}\textbf{Kategori} \\ \hline
\endfirsthead
\hline
\cellcolor[HTML]{D9EAF7}\textbf{Aspek} & \cellcolor[HTML]{D9EAF7}\textbf{Rerata Skor} & \cellcolor[HTML]{D9EAF7}\textbf{Nilai Maksimal} & \cellcolor[HTML]{D9EAF7}\textbf{Indeks} & \cellcolor[HTML]{D9EAF7}\textbf{Kategori} \\ \hline
\endhead
Kualitas Sistem & & & & \\ \hline
Kualitas Informasi & & & & \\ \hline
Kualitas Layanan & & & & \\ \hline
Kepuasan Pengguna & & & & \\ \hline
\textbf{Keseluruhan} & & & & \\ \hline
\end{longtable}

\textit{Indeks keseluruhan dapat dihitung dari rata-rata Rerata Skor seluruh pertanyaan dibagi Nilai Maksimal.}

\section{Masukan Terbuka dari Responden}

\begin{longtable}{|L{0.20}|L{0.42}|L{0.38}|}
\caption{Masukan Terbuka dari Responden}\\
\hline
\cellcolor[HTML]{D9EAF7}\textbf{Responden} & \cellcolor[HTML]{D9EAF7}\textbf{Masukan} & \cellcolor[HTML]{D9EAF7}\textbf{Tindak Lanjut} \\ \hline
\endfirsthead
\hline
\cellcolor[HTML]{D9EAF7}\textbf{Responden} & \cellcolor[HTML]{D9EAF7}\textbf{Masukan} & \cellcolor[HTML]{D9EAF7}\textbf{Tindak Lanjut} \\ \hline
\endhead
 & & \\ \hline
 & & \\ \hline
 & & \\ \hline
 & & \\ \hline
\end{longtable}

% Bab 8: Kesimpulan dan Tindak Lanjut
\chapter{Kesimpulan dan Tindak Lanjut}
\label{bab:kesimpulan-dan-tindak-lanjut}

\section{Ringkasan Hasil Pengujian Sementara (25 November 2026)}

\begin{enumerate}
  \item \textbf{Pengujian Otomatis (Unit Test \& Static Test) pada branch \texttt{test} (berbasis \texttt{main} \texttt{442bab4} + unit test \texttt{b7d46d9}):}
  \begin{itemize}
    \item \textbf{Frontend Unit Test (Jest):} 29/29 suite lolos (288/288 test), coverage 66,94\% baris.
    \item \textbf{Backend Unit Test (Vitest):} 2/2 file lolos (37/37 test).
    \item \textbf{Frontend Typecheck (\texttt{tsc --noEmit}):} Lolos, 0 error (defect \texttt{DF-01} telah ditutup setelah \texttt{npm ci}).
    \item \textbf{Frontend Lint (\texttt{expo lint}):} 0 error, 65 warning.
    \item \textbf{Backend Typecheck (\texttt{tsc}) \& Prisma (\texttt{prisma validate}):} Lolos, 0 error dan skema valid.
  \end{itemize}
  \item \textbf{Status Alpha Testing (51 Test Case):} Sebanyak \textbf{21 dari 51 TC (41,2\%)} telah terverifikasi lolos pada level logika, dataset 31 POI, dan API (\texttt{F03} sebagian, \texttt{F05} sebagian, \texttt{F07}, \texttt{F08}, \texttt{F10}, \texttt{F11}, \texttt{F12}). Sebanyak \textbf{30 TC lainnya} (yang memerlukan GPS nyata, kompas fisik, notifikasi lock screen, stopwatch layar, dan log proxy HTTPS) masih menunggu eksekusi lapangan pada perangkat \textbf{Samsung Galaxy A26} di Sudirman Street Food.
  \item \textbf{Status Beta Testing (UAT):} Belum dilaksanakan; instrumen tugas (\texttt{T1--T5}) dan kuesioner skala Likert (\texttt{KS1--KP3}) siap digunakan pada sesi UAT bersama responden.
\end{enumerate}

\section{Tindak Lanjut}

\begin{enumerate}
  \item Melaksanakan 30 test case lapangan (\texttt{F01}, \texttt{F02}, \texttt{F03-2/4}, \texttt{F04}, \texttt{F05-4}, \texttt{F06}, \texttt{F09}) di Sudirman Street Day and Night Market menggunakan Samsung Galaxy A26, lengkap dengan dokumentasi screenshot dan log proxy HTTPS.
  \item Melaksanakan Beta Testing (UAT) dengan minimal 2 responden di luar tim pengembang, lalu menghitung Indeks Penerimaan pada Tabel 7.3.
\end{enumerate}


% ===================== Lampiran & Referensi =====================
\appendix
% Lampiran dan Referensi
\chapter{Bukti Pengujian}
\label{lampiran:bukti-pengujian}

\textit{Lampirkan bukti pengujian: tautan video uji lapangan, screenshot per test case (diberi nomor sesuai TC ID), log pengujian (termasuk log proxy HTTPS untuk verifikasi F09-3), dan formulir kuesioner.}

\chapter*{Referensi}
\addcontentsline{toc}{chapter}{Referensi}

Menora, T., Primasari, C. H., Wibisono, Y. P., Sidhi, T. A. P., Setyohadi, D. B., \& Cininta, M. (2023). Implementasi pengujian alpha dan beta testing pada aplikasi Gamelan Virtual Reality. \textit{KONSTELASI: Konvergensi Teknologi dan Sistem Informasi}, 3(1), 48--60.


\end{document}
